%%%%%%%%%%%%%%%%%%%%%%%%%%%%%%%%%%%%%%%%%%%%%%%%%%%%%%%%%%%%%%%%%%%%%%%%
%% Congruent Dissections of the Regular Tetrahedron
%% A Computational Atlas and Algebraic Synthesis
%%
%% Author: Oleksiy Babanskyy
%%
%% Compile with:  pdflatex paper && pdflatex paper
%%                (two runs to resolve cross-references)
%%%%%%%%%%%%%%%%%%%%%%%%%%%%%%%%%%%%%%%%%%%%%%%%%%%%%%%%%%%%%%%%%%%%%%%%

\documentclass[11pt,a4paper]{article}

% ---------- Encoding & language ----------
\usepackage[T1]{fontenc}
\usepackage[utf8]{inputenc}
\usepackage[english]{babel}
\usepackage{lmodern}
\usepackage[protrusion=true,expansion=false]{microtype}

% ---------- Math ----------
\usepackage{amsmath,amssymb,amsthm}
\usepackage{aliascnt}
\usepackage{mathtools}
\usepackage{enumitem}

% ---------- Tables & floats ----------
\usepackage{booktabs}
\usepackage{array}
\usepackage{geometry}
\usepackage{longtable}
\usepackage{tabularx}
\usepackage{multirow}
\usepackage{float}
\usepackage{needspace}
\geometry{margin=2.5cm}

% ---------- Graphics & diagrams ----------
\usepackage{graphicx}
\usepackage{tikz}
\usetikzlibrary{arrows.meta, positioning, shapes.geometric, fit, calc}

% ---------- Code listings ----------
\usepackage{listings}
\lstset{
  basicstyle=\ttfamily\footnotesize,
  commentstyle=\color{gray!70!black}\itshape,
  keywordstyle=\color{blue!60!black}\bfseries,
  stringstyle=\color{red!50!black},
  breaklines=true,
  frame=single,
  rulecolor=\color{black!30},
  numbers=left,
  numberstyle=\tiny\color{gray},
  numbersep=6pt,
  showstringspaces=false,
  columns=fullflexible,
  keepspaces=true,
}

% ---------- Colours ----------
\usepackage{xcolor}
\definecolor{symblue}{RGB}{15,82,186}
\definecolor{lered}{RGB}{210,43,43}
\definecolor{axgreen}{RGB}{34,139,34}
\definecolor{parentgold}{RGB}{218,165,32}

% ---------- Hyperlinks (load late) ----------
\usepackage[colorlinks=true,
            linkcolor=blue!50!black,
            citecolor=blue!50!black,
            urlcolor=blue!70!black]{hyperref}
\usepackage[nameinlink,capitalise]{cleveref}
\Urlmuskip=0mu plus 2mu\relax
\setlength{\emergencystretch}{2em}

% ---------- Theorem environments ----------
\theoremstyle{plain}
\newtheorem{theorem}{Theorem}[section]
\newaliascnt{lemma}{theorem}
\newtheorem{lemma}[lemma]{Lemma}
\aliascntresetthe{lemma}
\newaliascnt{proposition}{theorem}
\newtheorem{proposition}[proposition]{Proposition}
\aliascntresetthe{proposition}
\newaliascnt{corollary}{theorem}
\newtheorem{corollary}[corollary]{Corollary}
\aliascntresetthe{corollary}
\newaliascnt{conditionaltheorem}{theorem}
\newtheorem{conditionaltheorem}[conditionaltheorem]{Conditional Theorem}
\aliascntresetthe{conditionaltheorem}
\newaliascnt{computationalcertificate}{theorem}
\newtheorem{computationalcertificate}[computationalcertificate]{Computational Certificate}
\aliascntresetthe{computationalcertificate}

\theoremstyle{definition}
\newaliascnt{definition}{theorem}
\newtheorem{definition}[definition]{Definition}
\aliascntresetthe{definition}
\newaliascnt{conjecture}{theorem}
\newtheorem{conjecture}[conjecture]{Conjecture}
\aliascntresetthe{conjecture}
\newaliascnt{question}{theorem}
\newtheorem{question}[question]{Question}
\aliascntresetthe{question}

\theoremstyle{remark}
\newaliascnt{remark}{theorem}
\newtheorem{remark}[remark]{Remark}
\aliascntresetthe{remark}
\newaliascnt{observation}{theorem}
\newtheorem{observation}[observation]{Observation}
\aliascntresetthe{observation}

% ---------- cleveref names ----------
\crefname{theorem}{theorem}{theorems}
\Crefname{theorem}{Theorem}{Theorems}
\crefname{lemma}{lemma}{lemmas}
\Crefname{lemma}{Lemma}{Lemmas}
\crefname{proposition}{proposition}{propositions}
\Crefname{proposition}{Proposition}{Propositions}
\crefname{corollary}{corollary}{corollaries}
\Crefname{corollary}{Corollary}{Corollaries}
\crefname{conditionaltheorem}{conditional theorem}{conditional theorems}
\Crefname{conditionaltheorem}{Conditional Theorem}{Conditional Theorems}
\crefname{computationalcertificate}{computational certificate}{computational certificates}
\Crefname{computationalcertificate}{Computational Certificate}{Computational Certificates}
\crefname{definition}{definition}{definitions}
\Crefname{definition}{Definition}{Definitions}
\crefname{conjecture}{conjecture}{conjectures}
\Crefname{conjecture}{Conjecture}{Conjectures}
\crefname{question}{question}{questions}
\Crefname{question}{Question}{Questions}
\crefname{remark}{remark}{remarks}
\Crefname{remark}{Remark}{Remarks}
\crefname{observation}{observation}{observations}
\Crefname{observation}{Observation}{Observations}

% ---------- Custom operators & macros ----------
\DeclareMathOperator{\Sym}{Sym}
\DeclareMathOperator{\Stab}{Stab}
\DeclareMathOperator{\Isom}{Isom}
\DeclareMathOperator{\Burn}{Burn}
\DeclareMathOperator{\Fix}{Fix}
\DeclareMathOperator{\Orb}{Orb}
\DeclareMathOperator{\Vol}{Vol}
\DeclareMathOperator{\aff}{aff}
\DeclareMathOperator{\conv}{conv}
\DeclareMathOperator{\rank}{rank}
\DeclareMathOperator{\BN}{BN}
\newcommand{\Td}{T_{\!d}}
\newcommand{\Q}{\mathbb{Q}}
\newcommand{\R}{\mathbb{R}}
\newcommand{\Z}{\mathbb{Z}}
\newcommand{\HH}{\mathbb{H}}
\newcommand{\calF}{\mathcal{F}}
\newcommand{\calO}{\mathcal{O}}
\newcommand{\Hroles}{H_{\mathrm{roles}}}
\newcommand{\Hslot}{H_{\mathrm{slot}}}
\newcommand{\Starp}{\mathrm{Star}^{+}}

% ---------- Metadata ----------
\hypersetup{
  pdftitle={Congruent Dissections of the Regular Tetrahedron},
  pdfauthor={Oleksiy Babanskyy},
  pdfsubject={Combinatorial geometry, Divisor Completeness, Tetrahedral rigidity},
  pdfkeywords={tetrahedron, dissection, Dehn invariant, Burnside ring, 24-cell, divisor completeness, transitivity conjecture}
}

\title{%
  Congruent Dissections of the Regular Tetrahedron:\\
  A Computational Atlas and Algebraic Synthesis}
\author{Oleksiy Babanskyy}
\date{}

%%%%%%%%%%%%%%%%%%%%%%%%%%%%%%%%%%%%%%%%%%%%%%%%%%%%%%%%%%%%%%%%%%%%%%%%
\begin{document}

\pagestyle{plain}

\maketitle

\begin{abstract}
We give an exact computational atlas for the eight known congruent dissections
of the regular tetrahedron $T$ into
$n \in \{1,2,3,4,6,8,12,24\}$ pieces.  Each representative is reconstructed
from canonical coordinates, its Dehn invariant is verified symbolically, and
its refinement relations are recorded.

Every atlas dissection is $G_{\mathrm{diss}}$-transitive, so Lagrange forces
$n \mid 24$ for the known transitive family; full $\Td$-invariance occurs only
for $n\in\{1,12,24\}$.  Burnside-ring, $24$-cell, and IFS-grammar models
realise every divisor and provide a tested screening stack for related
$\Td$-invariant orbit decompositions.

Beyond the atlas, we record two further scoped results.  For $n=16$ an
all-one-face case is excluded under two finite MILP hypotheses, automatic in
the Coxeter-pure instance.  For $n=5$, we prove a conditional obstruction under an explicit boundary-role
homogeneity hypothesis (BRH): five congruent convex pieces cannot partition the
regular tetrahedron when the same facets of the reference piece are exposed in
every copy.  The unrestricted case with different boundary roles remains open.
We also record exact real-area bookkeeping and a chart-level algebraic check;
these are not promoted to a general impossibility theorem.  The new $n=5$
computations use exact rational or symbolic arithmetic and are reproducible
from the repository scripts.
\end{abstract}

\smallskip
{\small
\noindent\textbf{Keywords:}
Regular tetrahedron, congruent dissection, Dehn invariant, tetrahedral symmetry group,
Burnside ring, 24-cell, $A_3$ Coxeter complex, Divisor Completeness conjecture,
Transitivity conjecture, face-to-face tiling.

\smallskip
\noindent\textbf{MSC\,2020:}
52B10, 20B25, 52C22, 51M20, 52B15, 05E18.
}

%%%%%%%%%%%%%%%%%%%%%%%%%%%%%%%%%%%%%%%%%%%%%%%%%%%%%%%%%%%%%%%%%%%%%%%%
\section{Introduction and summary of results}\label{sec:intro}

The dissection of the regular tetrahedron $T \subset \R^{3}$ into $n$ strictly
congruent pieces is a classical problem in combinatorial geometry.  Despite a
century of piecewise progress\,---\,from Hill's parallelepipeds
\cite{Hill1895} through the median-plane decompositions of Alexanderson and
Wetzel~\cite{AW1971}, the algorithmic refinements of Ong~\cite{Ong1991}, the
longest-edge bisection of Padr\'on--Plaza--Su\'arez~\cite{PPS2023}, the Dehn
invariant technology following Hadwiger~\cite{Hadwiger1951} and
Conway--Jones~\cite{ConwayJones1976}, and the symmetry-ring language of
Burnside~\cite{Burnside1911,tDieck1979}\,---\,no unified treatment has
appeared.  In particular the following natural conjecture has remained open:

\begin{conjecture}[Divisor Completeness]\label{conj:divisor}
$n$ admits a congruent dissection of the regular tetrahedron if and only if
$n \mid 24$.
\end{conjecture}

The purpose of the present paper is fourfold.

\medskip
\noindent\emph{First,} we provide a computational atlas of the
eight historically known dissections $n \in \{1,2,3,4,6,8,12,24\}$
(\Cref{sec:atlas}), with canonical Cartesian coordinates, exact squared-edge
spectra, and independently verified metric data.  We establish in particular
that \emph{every} piece of each of the eight historical representatives is a
tetrahedron (simplex with $f$-vector $(4,6,4)$) -- in contrast to a
persistent misconception in informal sources.  The atlas covers the
historical divisor representatives and does not aim at classifying arbitrary
convex Coxeter-pure partitions; see the companion
work~\cite{CoxeterPure2026} for the convex Coxeter-pure catalogue, which
includes pieces other than tetrahedra (e.g.\ an H-shape pentahedron at
$n=8$).

\noindent\emph{Second,} we determine the refinement poset on the eight
historical representatives
(\Cref{sec:poset}--\Cref{sec:refinement}), including the $S_4 \to S_8$
refinement bridging the Alexanderson--Wetzel and Padr\'on--Plaza--Su\'arez
traditions, the disproof of the conjectured $S_6 \to S_{12}$ refinement, and
the unique \emph{heterogeneous} bridge $S_2 = S_3 + S_6$ at cut parameter
$t = 2/3$.

\noindent\emph{Third,} we develop three independent algebraic obstruction
 frameworks for \cref{conj:divisor}: the Burnside ring $\Burn(\Td)$ with its
 $11$-class subgroup lattice (\Cref{sec:burnside}); the $24$-cell and binary
 tetrahedral correspondence $2T \leftrightarrow \Td \leftrightarrow \mathrm{flags}(T)$
 (\Cref{app:24cell}); and a formal grammar of affine contractions encoding all
 halving refinements (\Cref{sec:ifs}).  Together these realise every divisor
 $n \mid 24$ and, in the $G_{\mathrm{diss}}$-transitive case, force $n \mid 24$ by
 Lagrange's theorem applied to $G_{\mathrm{diss}}(D) \leq \Td$
 (\Cref{def:gdiss-transitive,cor:forward}).  The remaining open gap is precisely:

\begin{conjecture}[Transitivity Conjecture]\label{conj:transitivity}
Every congruent dissection of the regular tetrahedron is $G_{\mathrm{diss}}$-transitive
in the sense of \Cref{def:gdiss-transitive}.
\end{conjecture}

\smallskip

\noindent\emph{Relation to the Coxeter-pure classification.}  The present
atlas is complementary to the companion paper~\cite{CoxeterPure2026},
which classifies all convex Coxeter-pure (chamber-union) congruent
partitions of~$T$ into~$15$ $S_{4}$-orbit families.  The two scopes are
distinct: \cite{CoxeterPure2026} fixes the $A_{3}$ Coxeter complex as the
underlying combinatorics but allows arbitrary convex pieces (including
non-tetrahedral ones, e.g.\ a pentahedral H-shape at $n=8$), whereas the
present paper fixes the eight historical divisor representatives and
develops obstruction frameworks for non-Coxeter / non-transitive
dissections.  The naming crosswalk
$T_{i} \leftrightarrow S_{j}$ of~\cite{CoxeterPure2026}
is compatible with the atlas labels $S_{n}$ used here; see
\cref{ssec:atlas-table}.

\noindent\emph{Fourth,} in \Cref{sec:programme} we attack the reverse
direction via a structured programme whose two principal new results are:

\begin{conditionaltheorem}[$n=5$ boundary-role obstruction, preview]\label{thm:n5-intro}
Under the explicit boundary-role homogeneity condition BRH, no partition of
$T$ into $5$ full-dimensional congruent convex pieces exists.  BRH requires
one fixed set of reference facets to be exposed in every copy; it is not a
consequence of congruence alone.  The unrestricted mixed-role case remains
open.  See \Cref{thm:n5,sec:prog-n5} for the precise statement and proof.
\end{conditionaltheorem}

The $n=5$ contribution is therefore a conditional theorem plus exact finite
bookkeeping, rather than an unconditional residual reduction.  The corrected
real-area enumeration and the local NN identity are included to make the
remaining hypotheses visible; the historical phase-26 reduction scripts are
retained as audit material but are not used as proof evidence.

\begin{conditionaltheorem}[$n=16$ all-one-face, preview]\label{thm:n16-intro}
Under the medial-role and slot-separation hypotheses $(\Hroles) + (\Hslot)$,
no face-to-face convex congruent dissection of $T$ into $16$ pieces can
place exactly one boundary face of every piece on a single common $T$-face
(general convex instance: \cref{thm:n16-aof-conditional}; Coxeter-pure
instance: unconditional, \cref{thm:n16-aof-coxeter}).
See \Cref{sec:prog-n16} for the full statement, the list of rigidity
lemmas used, and the proof.
\end{conditionaltheorem}

\medskip

\noindent\textbf{Evidence levels.}
We adopt throughout the evidence-level tagging scheme \textsf{L0} (heuristic),
\textsf{L1} (partial evidence), \textsf{L1+} (strong partial evidence),
\textsf{L2} (conditional theorem), \textsf{L2+} (finite-check conditional
theorem), \textsf{L3} (unconditional theorem); see \Cref{sec:evidence-glossary}
for the precise definitions.  Every claim in this paper is tagged with one
of these six levels, and nothing else: hybrid tags such as \textsf{L2$^{-}$}
or \textsf{L1+}/\textsf{L2$^{-}$} are \emph{not} used.  The associated
computational artifacts (source code, test suites, JSON certificates) are
publicly available in the companion repository; see the \emph{Data and code
availability} statement at the end of \Cref{sec:intro}.

\medskip

\noindent\textbf{Organisation.}
\Cref{sec:atlas} presents the atlas; \Cref{sec:dehn}--\Cref{sec:cube} the
Dehn invariant verification, the refinement poset, transitivity plus the
Lagrange obstruction, and the comparison with the cube.
\Cref{sec:burnside}--\Cref{sec:ifs} develop the three algebraic obstruction
frameworks (Burnside ring, IFS grammar; the $24$-cell viewpoint is
deferred to \Cref{app:24cell}); \Cref{sec:nontrans} the
non-transitive impossibility analysis (obstruction stack under
$\Td$-invariance); \Cref{sec:programme} the rigidity and
configurational programme toward \cref{conj:transitivity}, including
\cref{thm:n5-intro,thm:n16-intro}.

\medskip

\noindent\textbf{Acknowledgements.}
Computational certificates were produced with Python \texttt{Fractions}
(exact rational arithmetic) and a custom MILP model built on
\texttt{PuLP}/\texttt{CBC}.

\medskip

\noindent\textbf{Data and code availability.}
The manuscript source, test suites, scripts, and JSON certificates released
with this paper are publicly available at
\url{https://github.com/aconsciousfractal/Congruent-Dissections-of-the-Regular-Tetrahedron}.
Executable artifacts are stored under \path{scripts/}, with deterministic
outputs under \path{results/}; these include
\path{results/S6_to_S12_edge_spectrum.json} and the active $n=5$ BRH,
area, and local-identity certificates.  The general-convex $n=16$ statement is explicitly conditional on
$(\Hroles)+(\Hslot)$; no separate public computational discharge of those two
hypotheses is claimed.  This release has no Zenodo DOI.

%%%%%%%%%%%%%%%%%%%%%%%%%%%%%%%%%%%%%%%%%%%%%%%%%%%%%%%%%%%%%%%%%%%%%%%%
\section{Preliminaries: four notions of symmetry}\label{sec:prelim}

The tetrahedral group $\Td \cong S_{4}$ acts isometrically on $T$ by permuting
its four vertices; this lifts to a linear action on $\R^{3}$ fixing the
centroid of $T$, and hence to an action on the power set of $T$ sending convex
sub-polyhedra to convex sub-polyhedra of the same congruence class.  In the
sequel we use four distinct notions of symmetry attached to a congruent
partition $D = \{P_{1},\ldots,P_{n}\}$ of $T$; they are logically independent
and have been repeatedly conflated in earlier literature.

\begin{definition}[Piece-type homogeneous space]\label{def:piece-type}
Fix a piece $P_{0} \in D$ and let
$\Stab_{\Td}(P_{0}) = \{ g \in \Td : g(P_{0}) = P_{0} \}$ be its pointwise
stabiliser under the $\Td$-action on convex sub-polyhedra of~$T$.  The
\emph{piece-type homogeneous space} of $P_{0}$ is the transitive $\Td$-set
$\Td / \Stab_{\Td}(P_{0})$.  Its cardinality equals the size of the
$\Td$-orbit $\Td \cdot P_{0} \subset 2^{T}$.
\end{definition}

\begin{definition}[Geometric symmetry group of a partition]\label{def:gdiss}
The \emph{geometric symmetry group} of $D$ is
\[
  G_{\mathrm{diss}}(D) \;:=\; \{ g \in \Td : g(D) = D \text{ as a set of pieces}\}
  \;\leq\; \Td,
\]
i.e.\ the subgroup of $\Td$ that permutes the pieces of $D$ among themselves.
\end{definition}

\begin{definition}[$\Td$-invariance]\label{def:td-invariance}
$D$ is \emph{$\Td$-invariant} if $G_{\mathrm{diss}}(D) = \Td$, i.e.\ every
$g \in \Td$ permutes the pieces of $D$.
\end{definition}

\begin{definition}[$G_{\mathrm{diss}}$-transitivity]\label{def:gdiss-transitive}
$D$ is \emph{$G_{\mathrm{diss}}$-transitive} if $G_{\mathrm{diss}}(D)$ (acting
by its inclusion $G_{\mathrm{diss}}(D) \leq \Td$) acts transitively on the set
of pieces of $D$: for every $P_{i}, P_{j} \in D$ there exists
$g \in G_{\mathrm{diss}}(D)$ with $g(P_{i}) = P_{j}$.  Equivalently, the $n$
pieces form a single orbit under the group of $\Td$-elements that preserve
$D$ set-wise.
\end{definition}

\begin{remark}[Logical relations]\label{rem:sym-relations}
The four notions are related as follows.
\begin{enumerate}
\item $\Td$-invariance (\Cref{def:td-invariance}) is the strongest
      condition and implies $G_{\mathrm{diss}}(D) = \Td$; in the atlas it
      holds only for $n \in \{1,12,24\}$ (\Cref{lem:geom-vs-alg-stab}).
\item $G_{\mathrm{diss}}$-transitivity (\Cref{def:gdiss-transitive}) is the
      notion that yields the Lagrange obstruction $n \mid 24$
      (\Cref{cor:forward}): orbit--stabiliser in $G_{\mathrm{diss}}(D)$ gives
      $n = [G_{\mathrm{diss}}(D) : \Stab_{0}]$, and
      $|G_{\mathrm{diss}}(D)|$ divides $|\Td| = 24$ by Lagrange applied to
      $G_{\mathrm{diss}}(D) \leq \Td$.
\item $\Td$-invariance $\Rightarrow$ $G_{\mathrm{diss}}$-transitivity
      (whenever $\Td$ itself acts transitively on $D$); the converse fails
      for $n \in \{2,3,4,6,8\}$, where $G_{\mathrm{diss}}(S_{n}) \subsetneq \Td$.
\item The \emph{piece-type} homogeneous space
      $\Td/\Stab_{\Td}(P_{0})$ of \Cref{def:piece-type} can have more
      elements than the partition $D$: an element
      $g \in \Td \setminus G_{\mathrm{diss}}(D)$ maps $P_{0}$ to a convex
      sub-polyhedron of $T$ that lies in a \emph{different} congruent
      partition.  Consequently $\Td$-type properties (e.g.\ the order
      $24/n$ quoted in \cref{tab:td-action}) are formal identities derived
      from~$|\Td|=24$ and $|D|=n$, not direct statements about pointwise
      stabilisers of specific realisations.
\end{enumerate}
Whenever the word \emph{transitive} appears unqualified in the sequel, it
refers to \Cref{def:gdiss-transitive}.
\end{remark}

%%%%%%%%%%%%%%%%%%%%%%%%%%%%%%%%%%%%%%%%%%%%%%%%%%%%%%%%%%%%%%%%%%%%%%%%
\section{The computational atlas \texorpdfstring{$\calF(T)$}{F(T)}}\label{sec:atlas}

Through exact symbolic computation we map each of the eight known dissections
into canonical Cartesian coordinates with full metric verification: $n$
congruent pieces, each of volume $\Vol(T)/n$, tiling $T$ without gap or
overlap.  Throughout this section the tetrahedron $T$ is embedded with
vertices
\begin{equation}\label{eq:atlas-coords}
A = (0,0,0),\; B = (1,0,0),\; C = \bigl(\tfrac12, \tfrac{\sqrt 3}{2}, 0\bigr),\;
D = \bigl(\tfrac12, \tfrac{\sqrt 3}{6}, \tfrac{\sqrt 6}{3}\bigr),
\end{equation}
so that $|AB| = |AC| = \cdots = |CD| = 1$. We write $M_{XY}$ for the midpoint
of edge $\overline{XY}$, $G_{XYZ}$ for the centroid of face $\triangle XYZ$,
and $G_T$ for the centroid of $T$.

\subsection{Historical atlas}\label{ssec:atlas-table}

\begin{table}[H]
\centering\small
\caption{The eight known congruent dissections of $T$.  Each piece has
$f$-vector $(4,6,4)$: all eight piece types are tetrahedra.  Squared edge
spectrum $E^{2}$ is normalised to the parent edge length $|AB| = 1$.}
\label{tab:atlas}
\begin{tabular}{@{}ccccl@{}}
\toprule
$n$ & Piece name & $f$-vector & $E^{2}$ (squared edges) & Construction \\
\midrule
$1$  & Regular tetrahedron & $(4,6,4)$ & $[1,1,1,1,1,1]$                   & Identity \\
$2$  & Median wedge        & $(4,6,4)$ & $[\tfrac14,\tfrac34,\tfrac34,1,1,1]$ & 1 median plane \\
$3$  & Axial wedge         & $(4,6,4)$ & $[\tfrac13,\tfrac13,\tfrac23,1,1,1]$ & 3 axial planes \\
$4$  & Quadrisection       & $(4,6,4)$ & $[\tfrac14,\tfrac14,\tfrac12,\tfrac34,\tfrac34,1]$ & 2 median planes \\
$6$  & Hexasection         & $(4,6,4)$ & $[\tfrac{1}{12},\tfrac14,\tfrac13,\tfrac23,\tfrac34,1]$ & 3 axial $+$ 3 median \\
$8$  & $8T$-LE piece       & $(4,6,4)$ & $[\tfrac14,\tfrac14,\tfrac14,\tfrac14,\tfrac12,\tfrac34]$ & longest-edge iteration \\
$12$ & $A_{4}$ domain      & $(4,6,4)$ & $[\tfrac{1}{24},\tfrac13,\tfrac13,\tfrac38,\tfrac38,1]$ & $A_{4}$ fundamental domain \\
$24$ & Orthoscheme         & $(4,6,4)$ & $[\tfrac{1}{24},\tfrac{1}{12},\tfrac18,\tfrac14,\tfrac13,\tfrac38]$ & 6 median planes \\
\bottomrule
\end{tabular}
\end{table}

\noindent\emph{Remark on the $n=12$ entry.}  By ``$A_{4}$ fundamental
domain'' we mean a fundamental domain for the action of the alternating
group $A_{4}$, the orientation-preserving subgroup of
$\Td \cong S_{4}$, on~$T$: any such domain is a $12$-piece dissection
piece, concretely realised as $\conv(A,G_{ABC},G_{T},B)$
(\cref{tab:atlas-reps}), i.e.~the convex hull of a vertex, a face
barycentre, the $T$-centroid, and an edge-adjacent vertex.

\begin{table}[H]
\centering\small
\caption{Canonical representative pieces in the embedding~\eqref{eq:atlas-coords}.}
\label{tab:atlas-reps}
\begin{tabular}{@{}cl@{}}
\toprule
$n$ & Representative $S_{n}$ \\
\midrule
$2$  & $\conv(A,B,C,M_{CD})$ \\
$3$  & $\conv(D,G_{ABC},A,B)$ \\
$4$  & $\conv(A,M_{AB},C,M_{CD})$ \\
$6$  & $\conv(D,G_{ABC},A,M_{AB})$ \\
$8$  & $\conv(A,M_{AB},M_{AC},M_{CD})$ \\
$12$ & $\conv(A,G_{ABC},G_{T},B)$ \\
$24$ & $\conv(A,M_{AB},G_{ABC},G_{T})$ \\
\bottomrule
\end{tabular}
\end{table}

\begin{theorem}[Tetrahedrality of the eight historical representatives]\label{thm:univ-tet}
Every piece of each of the eight historical congruent dissection
representatives $\{S_{n}\}_{n \in \mathrm{div}(24)}$ listed in
\cref{tab:atlas-reps} is a tetrahedron.
\end{theorem}

\begin{proof}
Direct verification from \cref{tab:atlas-reps}: each representative is the
convex hull of exactly four affinely independent points.  The informal
description of $S_{3}$ as a ``wedge'' pentahedron in certain forum sources is
a notational artefact and not the actual piece geometry; similarly, the
$S_{12}$ piece obtained by merging two adjacent orthoschemes reduces to a
tetrahedron because the merged midpoint $M_{AB}$ lies on edge $\overline{AB}$
and drops out of the convex hull.
\end{proof}

\begin{remark}[Scope]\label{rem:univ-tet-scope}
\cref{thm:univ-tet} is a statement about the eight specific historical
representatives, not about all Coxeter-pure congruent partitions of~$T$.
The companion classification~\cite{CoxeterPure2026} exhibits convex
Coxeter-pure pieces which are not tetrahedra; in particular at $n=8$ the
H-shape family~$\mathcal{F}_{11}$ has $f$-vector $(6,9,5)$ and is a
pentahedron.  The present atlas fixes the eight classical
representatives; the general convex Coxeter-pure case is
addressed in~\cite{CoxeterPure2026}.
\end{remark}

%%%%%%%%%%%%%%%%%%%%%%%%%%%%%%%%%%%%%%%%%%%%%%%%%%%%%%%%%%%%%%%%%%%%%%%%
\section{Dehn invariants of the atlas pieces}\label{sec:dehn}

The Dehn invariant $D(P) \in \R \otimes_{\Z} (\R/\pi\Q)$ is additive under
scissors decomposition (Sydler~\cite{Sydler1965}; see also
Dupont~\cite{Dupont2001} for the modern group-homological
formulation), and the regular tetrahedron carries the irrational
value $D(T) = 6 \otimes \alpha$ where
$\alpha \coloneqq \arccos(1/3) \approx 70.529^{\circ}$ is the dihedral angle
along every edge.  By additivity a congruent $n$-dissection $T = \bigcup S_{n}$
requires
\begin{equation}\label{eq:dehn-scaling}
  D(S_{n}) = \frac{1}{n} D(T) = \frac{6}{n}\otimes\alpha .
\end{equation}

Computation of the piece parameters reveals a conjugate dihedral angle
$\beta \coloneqq \arccos(\sqrt{6}/3)$.  By the trigonometric double-angle
identity one finds $\cos(2\beta) = 2 \cdot (6/9) - 1 = 1/3$.  Since both
$2\beta$ and $\alpha$ lie in $(0,\pi)$, this gives the exact identity
\begin{equation}\label{eq:2beta-alpha}
  2\beta = \alpha.
\end{equation}

This local angle identity is not a classification of all individual
dihedral angles of the atlas pieces.  The total Dehn tensors below follow
from additivity and congruence, independently of such a classification.

\begin{corollary}[Dehn values of the atlas pieces]\label{thm:dehn-values}
Applying~\eqref{eq:dehn-scaling} to the eight congruent atlas dissections
gives $D(S_1)=6\otimes\alpha$ and the following values:
\begin{center}\small
\begin{tabular}{@{}ccccccc@{}}
\toprule
$n$ & $2$ & $3$ & $4$ & $6$ & $8$ & $12$ \\[2pt]
$D(S_{n})$ & $3\otimes\alpha$ & $2\otimes\alpha$ & $\tfrac{3}{2}\otimes\alpha$ &
$1\otimes\alpha$ & $\tfrac34\otimes\alpha$ & $\tfrac12\otimes\alpha$ \\
\bottomrule
\end{tabular}
\quad
\begin{tabular}{@{}cc@{}}
\toprule
$n$ & $24$ \\[2pt]
$D(S_{n})$ & $\tfrac14\otimes\alpha$ \\
\bottomrule
\end{tabular}
\end{center}
\end{corollary}

\begin{corollary}[No cube-equivalent piece]\label{cor:nocube}
No piece $S_{n}$ of any known congruent dissection of $T$ has $D(S_{n}) = 0$;
consequently no such piece is scissors-congruent to a cube.
\end{corollary}

%%%%%%%%%%%%%%%%%%%%%%%%%%%%%%%%%%%%%%%%%%%%%%%%%%%%%%%%%%%%%%%%%%%%%%%%
\section{The structural lineage \texorpdfstring{$S_{4} \to S_{8}$}{S4->S8}}
\label{sec:poset}

The $8T$-LE construction of Padr\'on, Plaza and Su\'arez~\cite{PPS2023} is
visually unlike the global median-plane decompositions of
Alexanderson--Wetzel~\cite{AW1971}: it is generated by iterated
\emph{longest-edge bisection}, a local algorithmic procedure.  Canonical
coordinates however reveal a clean lineage.

\begin{observation}[Canonical lineage]\label{obs:s4-s8}
The $8T$-LE dissection descends from the $S_{4}$ quadrisection as follows.
\begin{enumerate}
  \item Bisecting the longest edge of $T$ (any edge, by symmetry) yields $S_{2}$.
  \item Bisecting the longest edge of $S_{2}$ introduces a strictly orthogonal
        plane (perpendicular to the first median plane), yielding $S_{4}$.
  \item The longest edge of $S_{4}$ is the unique remaining original $T$-edge
        of length $1$; bisecting it divides $S_{4}$ into two copies of $S_{8}$.
\end{enumerate}
\end{observation}

The refinement graph over halving cuts is summarised in
\Cref{fig:refinement}.  The axial and symmetry-plane pipelines are
geometrically \emph{incompatible} at the $S_{6} \to S_{12}$ junction: the
$S_{6}$ representative $\conv(D,G_{ABC},A,M_{AB})$ uses the apex $D$, while
$S_{12} = \conv(A,G_{ABC},G_{T},B)$ uses the tetrahedron centroid $G_{T}$ and
a fresh vertex $B$.  The two families thus form disjoint components
connected only through the heterogeneous bridge $S_{2} = S_{3} + S_{6}$
described in~\Cref{sec:refinement}.

\begin{figure}[H]
\centering
\begin{tikzpicture}[
  >=Stealth,
  node/.style={rectangle, rounded corners=2pt, draw, minimum width=30mm,
               minimum height=7mm, align=center, font=\small\bfseries,
               thick},
  sym/.style ={node, fill=symblue!85,    text=white, draw=symblue!40!black},
  ax/.style  ={node, fill=axgreen!80,    text=white, draw=axgreen!40!black},
  le/.style  ={node, fill=lered!85,      text=white, draw=lered!40!black},
  parent/.style={node, fill=parentgold!80, text=black, draw=black, thick,
                 minimum width=36mm},
  arr/.style ={->, thick},
  darr/.style={->, thick, dashed},
  xarr/.style={->, thick, dotted, red!70!black},
  every node/.append style={font=\small}
]
\node [parent] (T)          at (0, 5)   {$T$: regular tetrahedron};

% Median pipeline (left)
\node [sym]    (S2)         at (-5, 3)  {$S_2$ median};
\node [sym]    (S4)         at (-5, 1)  {$S_4$ quadrisection};
\node [sym]    (S12)        at (-5,-1)  {$S_{12}$ dodecasection};
\node [sym]    (S24)        at (-5,-3)  {$S_{24}$ orthoscheme};

% Axial pipeline (right)
\node [ax]     (S3)         at ( 5, 3)  {$S_3$ axial trisection};
\node [ax]     (S6)         at ( 5, 1)  {$S_6$ hexasection};

% LE pipeline
\node [le]     (S8)         at (-1.2,-1){$S_8$ $8T$-LE};

\draw [arr] (T)  -- node[left,midway,font=\tiny]{1 median}            (S2);
\draw [arr] (S2) -- node[left,midway,font=\tiny]{$\perp$ median}      (S4);
\draw [arr] (S4) -- node[left,midway,font=\tiny]{$A_4$ domain}        (S12);
\draw [arr] (S12)-- node[left,midway,font=\tiny]{final 2 planes}      (S24);

\draw [arr] (T)  -- node[right,midway,font=\tiny]{3 axial}            (S3);
\draw [arr] (S3) -- node[right,midway,font=\tiny]{base bisect}        (S6);

\draw [arr] (S4) -- node[above,midway,font=\tiny]{LE bisect}         (S8);

\draw [darr] (S2) to[bend right=25] node[above,font=\tiny]
                     {$S_2 = S_3 + S_6$} (S3);
\draw [darr] (S2) to[bend right=35]                           (S6);

\draw [xarr] (S6) to[bend right=12] node[below,font=\tiny]{\textbf{not geometric}} (S12);
\end{tikzpicture}
\caption{The refinement poset of the eight known dissections. Solid arrows
are verified halving refinements; dashed arrows record the unique
heterogeneous bridge $S_2 = S_3 + S_6$ (at cut parameter $t = 2/3$, see
\Cref{sec:refinement}); the dotted red arrow marks the geometrically
\emph{disproved} refinement $S_6 \to S_{12}$.  Blue = symmetry pipeline,
green = axial pipeline, red = longest-edge pipeline.}
\label{fig:refinement}
\end{figure}

\begin{theorem}[Disproof of $S_{6} \to S_{12}$]\label{thm:s6-s12-fail}
There is no congruent halving of the $S_{6}$ piece into two tetrahedra
each congruent to $S_{12}$.
\end{theorem}

\begin{proof}
The argument is by the squared-edge multiset, a congruence invariant:
if two tetrahedra $P, Q$ are congruent then the unordered multisets
$E^{2}(P), E^{2}(Q)$ of squared edge lengths coincide, and in particular
$\min E^{2}(P) = \min E^{2}(Q)$.  From \cref{tab:atlas} we read
\begin{equation}\label{eq:s12-spectrum}
  E^{2}(S_{12}) = \Bigl\{\tfrac{1}{24},\,\tfrac13,\,\tfrac13,\,\tfrac38,\,\tfrac38,\,1\Bigr\},
  \qquad \min E^{2}(S_{12}) = \tfrac{1}{24}.
\end{equation}

\smallskip\noindent\emph{Shape of a congruent halving.}
A plane cut of a tetrahedron $P$ produces two tetrahedra if and only if
the cut plane contains exactly one edge of $P$ and meets the opposite
edge at a single interior point (any cut meeting three edges produces a
tetrahedron and a triangular prism, hence cannot yield two tetrahedra).
If in addition the two halves are congruent -- and hence of equal
volume $\tfrac12\Vol(P)$ -- the interior intersection point must be the
midpoint of the opposite edge.  Applying this to
$S_{6} = \conv(D, G_{ABC}, A, M_{AB})$ leaves exactly three opposite-edge
pairs and hence six candidate half-tetrahedra.

\smallskip\noindent\emph{Exact enumeration in the embedding~\eqref{eq:atlas-coords}.}
A direct computation in rational arithmetic (certificate
\path{results/S6_to_S12_edge_spectrum.json} in the public companion
repository; see the \emph{Data and code availability} statement in
\Cref{sec:intro})
gives the
following squared-edge multisets for the six candidate halves, indexed
by the edge $e$ preserved by the cut plane and the vertex $v$ of
$\{D, G_{ABC}, A, M_{AB}\} \setminus e$ retained in that half:
\begin{align*}
  e = \{D,G_{ABC}\}:\ &\ \bigl\{\tfrac{1}{16}, \tfrac{7}{48}, \tfrac13, \tfrac23, \tfrac{13}{16}, 1\bigr\},\
  \bigl\{\tfrac{1}{16}, \tfrac{1}{12}, \tfrac{7}{48}, \tfrac23, \tfrac34, \tfrac{13}{16}\bigr\}; \\
  e = \{D,A\}:\        &\ \bigl\{\tfrac{1}{48}, \tfrac{13}{48}, \tfrac13, \tfrac23, \tfrac{11}{16}, 1\bigr\},\
  \bigl\{\tfrac{1}{48}, \tfrac14, \tfrac{13}{48}, \tfrac{11}{16}, \tfrac34, 1\bigr\}; \\
  e = \{D,M_{AB}\}:\   &\ \bigl\{\tfrac{1}{12}, \tfrac{1}{12}, \tfrac{1}{12}, \tfrac23, \tfrac34, \tfrac34\bigr\},\
  \bigl\{\tfrac{1}{12}, \tfrac{1}{12}, \tfrac14, \tfrac34, \tfrac34, 1\bigr\}.
\end{align*}
The set of minima attained across the six halves is
$\{\tfrac{1}{48}, \tfrac{1}{16}, \tfrac{1}{12}\}$.  By
\eqref{eq:s12-spectrum}, $\min E^{2}(S_{12}) = \tfrac{1}{24}$ does not
belong to this set, so no half has $S_{12}$'s squared-edge multiset.
Hence no congruent halving of $S_{6}$ produces two copies of $S_{12}$.
\end{proof}

%%%%%%%%%%%%%%%%%%%%%%%%%%%%%%%%%%%%%%%%%%%%%%%%%%%%%%%%%%%%%%%%%%%%%%%%
\section{The refinement graph}\label{sec:refinement}

\subsection{Egyptian fractions over \texorpdfstring{$\mathrm{div}(24)$}{div(24)}}

We exhaustively test single-plane congruent cuts
$S_{a} \to S_{b} \cup S_{c}$ with $a, b, c \in \mathrm{div}(24)$,
$b \leq c$, $b > a$ and volume-consistent Egyptian fraction
$1/a = 1/b + 1/c$; the target is geometric realisability
(edge-spectrum congruence of the two halves).

\begin{table}[H]
\centering\small
\caption{Egyptian fraction audit.  Six of the eleven admissible identities
are geometrically realisable.}
\label{tab:egyptian}
\begin{tabular}{@{}lcl@{}}
\toprule
Identity & Geometric? & Note \\
\midrule
$1 = \tfrac12 + \tfrac12$               & \checkmark & Median bisection \\
$\tfrac12 = \tfrac13 + \tfrac16$        & \checkmark (new)  & Cross-pipeline bridge \\
$\tfrac12 = \tfrac14 + \tfrac14$        & \checkmark & Orthogonal median \\
$\tfrac13 = \tfrac14 + \tfrac{1}{12}$   & $\times$   & \\
$\tfrac13 = \tfrac16 + \tfrac16$        & \checkmark & Base median bisection \\
$\tfrac14 = \tfrac16 + \tfrac{1}{12}$   & $\times$   & \\
$\tfrac14 = \tfrac18 + \tfrac18$        & \checkmark & $8T$-LE bisection \\
$\tfrac16 = \tfrac18 + \tfrac{1}{24}$   & $\times$   & \\
$\tfrac16 = \tfrac{1}{12} + \tfrac{1}{12}$ & $\times$   & $S_6 \to S_{12}$ disproved \\
$\tfrac18 = \tfrac{1}{12} + \tfrac{1}{24}$ & $\times$   & volumes match, shapes do not \\
$\tfrac{1}{12} = \tfrac{1}{24} + \tfrac{1}{24}$ & \checkmark & orthoscheme split \\
\bottomrule
\end{tabular}
\end{table}

\subsection{The heterogeneous bridge \texorpdfstring{$S_{2} = S_{3} + S_{6}$}{S2 = S3 + S6}}

\begin{theorem}[Cross-pipeline bridge]\label{thm:s2-bridge}
The single-plane cut of the $S_{2}$ representative
$S_{2} = \conv(A, B, C, M_{CD})$ at cut parameter $t = 2/3$ along a cut-edge
produces one piece with normalised squared edges $(1,1,2,3,3,3)$
(matching $S_{3}$) and one piece with normalised squared edges
$(1,3,4,8,9,12)$ (matching $S_{6}$).  This is the unique heterogeneous
refinement bridge in the catalogue.
\end{theorem}

The two concrete realisations are related by swapping the fixed edge; in
both, $\det(\alpha) \cdot \det(\beta) = -1$, so one refinement map is
orientation-preserving and one is orientation-reversing.

\subsection{Connected components}

\begin{proposition}[Component structure]\label{prop:components}
Under homogeneous ($1{:}1$) halving refinements the refinement graph has
\emph{three} disconnected components:
\begin{center}
\begin{tabular}{@{}cll@{}}
\toprule
Component & Nodes                               & Pipeline \\
\midrule
1 & $\{S_{1}, S_{2}, S_{4}, S_{8}\}$            & Symmetry $+$ LE \\
2 & $\{S_{3}, S_{6}\}$                          & Axial \\
3 & $\{S_{12}, S_{24}\}$                        & $A_{4}$ / Orthoscheme \\
\bottomrule
\end{tabular}
\end{center}
The heterogeneous bridge of \cref{thm:s2-bridge} connects components $1$
and $2$; components $2$ and $3$ are geometrically disjoint.
\end{proposition}

%%%%%%%%%%%%%%%%%%%%%%%%%%%%%%%%%%%%%%%%%%%%%%%%%%%%%%%%%%%%%%%%%%%%%%%%
\section{\texorpdfstring{$\Td$}{Td}-transitivity and the Lagrange obstruction}
\label{sec:transitivity}

\subsection{All eight known dissections are $G_{\mathrm{diss}}$-transitive}
\label{ssec:transitivity}

\begin{computationalcertificate}[$G_{\mathrm{diss}}$-transitivity of the atlas]\label{thm:transitivity}
For every $n \in \{1,2,3,4,6,8,12,24\}$ the atlas dissection $S_{n}$ is
$G_{\mathrm{diss}}$-transitive in the sense of
\Cref{def:gdiss-transitive}: the geometric symmetry group
$G_{\mathrm{diss}}(S_{n}) \leq \Td$ of \Cref{def:gdiss} acts transitively on
the $n$ pieces of $S_{n}$, with orbit--stabiliser identity
$|G_{\mathrm{diss}}(S_{n})| = n \cdot |\Stab_{0}|$
(\Cref{lem:geom-vs-alg-stab}).  The numerical data of this certificate are
tabulated in \cref{tab:td-action,tab:gdiss}.
\end{computationalcertificate}

\begin{remark}[Reading \cref{tab:td-action}]\label{rem:type-vs-invariance}
The column $|\Stab|$ of \cref{tab:td-action} records the formal identity
$|\Stab| = 24/n$; for $n \in \{1,12,24\}$ it coincides with the
pointwise stabiliser of a piece in $\Td$ (since $G_{\mathrm{diss}} = \Td$),
while for $n \in \{2,3,4,6,8\}$ the pointwise $\Td$-stabiliser of a
specific realisation is strictly smaller and the correct orbit--stabiliser
decomposition takes place inside $G_{\mathrm{diss}}(S_{n}) \subsetneq \Td$
(\Cref{lem:geom-vs-alg-stab}).  The value $24/n$ nonetheless correctly
controls the Lagrange obstruction via $|G_{\mathrm{diss}}| \mid 24$.
\end{remark}

\begin{table}[H]
\centering\small
\caption{The $\Td$-action on each of the eight known dissections.}
\label{tab:td-action}
\begin{tabular}{@{}c c c l@{}}
\toprule
$n$ & Orbit size & $|\Stab|$ & Stabiliser type \\
\midrule
$1$  & $1$  & $24$ & $\Td$ \\
$2$  & $2$  & $12$ & $A_{4}$ \\
$3$  & $3$  & $8$  & $D_{2d}$ \\
$4$  & $4$  & $6$  & $S_{3}$ \\
$6$  & $6$  & $4$  & $V_{4}$ (Klein four-group) \\
$8$  & $8$  & $3$  & $\Z_{3}$ \\
$12$ & $12$ & $2$  & $\Z_{2}$ \\
$24$ & $24$ & $1$  & $\{e\}$ \\
\bottomrule
\end{tabular}
\end{table}

\begin{remark}[On the $n=8$ anomaly]\label{rem:n8-anomaly}
Although the $8T$-LE dissection is generated by an iterative algorithmic
procedure rather than as a subgroup orbit, $\Td$ nevertheless acts
transitively on its $8$ pieces, with stabiliser $\Z_{3}$ generated by the
$120^{\circ}$ rotation about a body diagonal.  In particular $8 \mid 24$ is
\emph{structural}, not accidental.
\end{remark}

\subsection{The Lagrange obstruction}\label{ssec:lagrange}

The immediate consequence of \cref{thm:transitivity} is:

\begin{corollary}[Divisor realisation and Lagrange obstruction]\label{cor:forward}
If $n \mid 24$, the regular tetrahedron admits a congruent dissection into
$n$ pieces.  Moreover, for any $G_{\mathrm{diss}}$-transitive congruent
dissection of $T$ (\Cref{def:gdiss-transitive}) necessarily $n \mid 24$.
\end{corollary}

Both statements follow by construction (\cref{tab:atlas}) and by applying
Lagrange's theorem to the chain $\Stab_{0} \leq G_{\mathrm{diss}}(D) \leq \Td$:
$n = [G_{\mathrm{diss}}(D) : \Stab_{0}]$ divides $|G_{\mathrm{diss}}(D)|$,
which in turn divides $|\Td| = 24$.  The full
divisor-completeness conjecture thus reduces to ruling out congruent
dissections that fail $G_{\mathrm{diss}}$-transitivity -- the subject of
\Cref{sec:nontrans,sec:programme}.

%%%%%%%%%%%%%%%%%%%%%%%%%%%%%%%%%%%%%%%%%%%%%%%%%%%%%%%%%%%%%%%%%%%%%%%%
\section{Comparison with the cube}\label{sec:cube}

To probe whether \cref{conj:divisor} generalises to other Platonic solids
we analyse the cube (symmetry group $O_{h}$, $|O_{h}| = 48$).

\begin{proposition}[Slab argument]\label{prop:slab}
The unit cube $[0,1]^{3}$ admits a congruent dissection into $n$ pieces for
\emph{every} $n \geq 1$: cut by planes $z = k/n$, $k = 1,\ldots,n-1$, to
produce $n$ congruent rectangular slabs $1 \times 1 \times (1/n)$, all
related by translation.
\end{proposition}

The mechanism of \cref{prop:slab} requires \emph{parallel faces} on the
parent body.  Among all Platonic solids the tetrahedron is the unique one
with \emph{zero} parallel face-pairs (\cref{tab:platonic}), so the slab
argument cannot be transplanted to it.

\begin{table}[H]
\centering\small
\caption{Parallel-face count and slab cuts among the Platonic solids.}
\label{tab:platonic}
\begin{tabular}{@{}lcccc@{}}
\toprule
Solid & Group & $|G|$ & Parallel face-pairs & Slab cuts? \\
\midrule
Tetrahedron  & $\Td$   & $24$  & $\mathbf{0}$ & \textbf{No} \\
Cube         & $O_{h}$ & $48$  & $3$          & Yes \\
Octahedron   & $O_{h}$ & $48$  & $4$          & Yes \\
Dodecahedron & $I_{h}$ & $120$ & $6$          & Yes \\
Icosahedron  & $I_{h}$ & $120$ & $6$          & Yes \\
\bottomrule
\end{tabular}
\end{table}

A refined question remains valid on the cube: for which $n$ does the cube
admit an \emph{$O_{h}$-transitive} congruent dissection?  The answer is
exactly $n \mid 48$, with explicit $O_{h}$-transitive constructions for all
ten divisors of $48$ (\cref{tab:oh-div}).

\begin{table}[H]
\centering\small
\caption{$O_{h}$-transitive congruent dissections of the cube.}
\label{tab:oh-div}
\begin{tabular}{@{}ccl@{}}
\toprule
$n$ & $|\Stab|$ & Construction \\
\midrule
$1$  & $48$ & Trivial \\
$2$  & $24$ & Mirror-plane bisection \\
$3$  & $16$ & Double-wedge (bicapped square) \\
$4$  & $12$ & Two orthogonal mirror planes \\
$6$  & $8$  & Six face pyramids \\
$8$  & $6$  & $2 \times 2 \times 2$ cubelets \\
$12$ & $4$  & Bisection of six pyramids \\
$16$ & $3$  & Grid subdivisions \\
$24$ & $2$  & $O$-orbit orthoscheme \\
$48$ & $1$  & $O_{h}$-orbit orthoscheme \\
\bottomrule
\end{tabular}
\end{table}

\begin{corollary}[Failure of divisor completeness for the cube]\label{cor:cube-fail}
Divisor completeness does not hold for the cube: all $n \geq 1$ admit
non-transitive slab dissections.  The emerging dichotomy is:
\begin{center}\small
\begin{tabular}{@{}lcc@{}}
\toprule
 & Tetrahedron ($\Td$) & Cube ($O_{h}$) \\
\midrule
$G$-transitive: $n \mid |G|$ sufficient? & \checkmark\ (8 divisors) & \checkmark\ (10 divisors) \\
Non-$G$-transitive constructions         & none known              & all $n$ (slabs) \\
Parallel faces                           & $0$                    & $3$ \\
\bottomrule
\end{tabular}
\end{center}
\end{corollary}

Hence \cref{conj:divisor} is \emph{not} a generic polyhedral phenomenon: it
is specifically the absence of parallel faces on the regular tetrahedron
that prevents non-transitive constructions and motivates
\cref{conj:transitivity}.

\begin{remark}[Arithmetic note]
The aggregate divisor sums do not transfer:
\[
  \frac{\sigma(24)}{24} = \frac{60}{24} = \frac{5}{2},
  \qquad
  \frac{\sigma(48)}{48} = \frac{124}{48} = \frac{31}{12}.
\]
\end{remark}

%%%%%%%%%%%%%%%%%%%%%%%%%%%%%%%%%%%%%%%%%%%%%%%%%%%%%%%%%%%%%%%%%%%%%%%%
\section{The Burnside ring invariant}\label{sec:burnside}

The Burnside ring $\Burn(\Td)$~\cite{Burnside1911,tDieck1979} provides a
representation-theoretic framework for characterising transitive $G$-sets
and detecting obstructions to non-transitive dissections.

\subsection{The subgroup lattice of \texorpdfstring{$\Td \cong S_{4}$}{Td=S4}}

The tetrahedral group $\Td \cong S_{4}$ has \emph{$11$ conjugacy classes of
subgroups} and $30$ subgroups in total (\cref{tab:s4-subgroups}).

\begin{table}[H]
\centering\small
\caption{The $11$ conjugacy classes of subgroups of $S_{4}$.}
\label{tab:s4-subgroups}
\begin{tabular}{@{}c l c c c@{}}
\toprule
Class & Subgroup $H$ & $|H|$ & $[S_{4}:H]$ & \# conjugates \\
\midrule
$0$  & $\{e\}$                                & $1$  & $24$ & $1$ \\
$1$  & $\Z_{2}^{t}$ (transposition)            & $2$  & $12$ & $6$ \\
$2$  & $\Z_{2}^{d}$ (double transposition)     & $2$  & $12$ & $3$ \\
$3$  & $\Z_{3}$                                & $3$  & $8$  & $4$ \\
$4$  & $V_{4}$ (non-normal Klein-4)            & $4$  & $6$  & $3$ \\
$5$  & $V_{4}^{n}$ (normal Klein-4)            & $4$  & $6$  & $1$ \\
$6$  & $\Z_{4}$                                & $4$  & $6$  & $3$ \\
$7$  & $S_{3}$                                 & $6$  & $4$  & $4$ \\
$8$  & $D_{4}$                                 & $8$  & $3$  & $3$ \\
$9$  & $A_{4}$                                 & $12$ & $2$  & $1$ \\
$10$ & $S_{4}$                                 & $24$ & $1$  & $1$ \\
\bottomrule
\end{tabular}
\end{table}

The valid orbit sizes for any $\Td$-action are thus
$\{1,2,3,4,6,8,12,24\}$ -- precisely the eight divisors of $24$.

\subsection{Equivariant Euler characteristic}

The barycentric subdivision of $T$ yields a $\Td$-equivariant
CW-complex.  A direct count gives
\begin{equation}\label{eq:chi-T}
  \chi(T) = 15 - 50 + 60 - 24 = 1,
\end{equation}
consistent with $T$ being a contractible $3$-ball.  Lifting to the Burnside
ring:
\begin{computationalcertificate}[Equivariant Euler class]\label{prop:chi-T}
The equivariant Euler characteristic $\chi_{\Td}(T)$ of $T$, viewed as an
equivariant CW-complex, equals the class of the point:
$\chi_{\Td}(T) = [\Td/\Td] \in \Burn(\Td)$, with character vector
$(1,1,\ldots,1) \in \Z^{11}$.
\end{computationalcertificate}

\subsection{The mark table}

The mark table $M[i][j] = |\Fix_{H_{j}}(\Td/H_{i})|$ encodes the fixed-point
structure of every transitive $\Td$-set.  Its computed value is given in
\Cref{tab:mark-full}; it has full rank $11$, confirming that $\Burn(\Td)$ is
the free abelian group of rank $11$ on the transitive $\Td$-sets.

\begin{table}[H]
\centering\footnotesize
\caption{Mark table of $\Td \cong S_{4}$.  Rows indexed by transitive $\Td$-sets
$\Td/H_{i}$ (top-to-bottom: subgroup order ascending); columns indexed by
subgroups $H_{j}$ (left-to-right: same ordering).}
\label{tab:mark-full}
\renewcommand{\arraystretch}{1.05}
\setlength{\tabcolsep}{3pt}
\begin{tabular}{@{} l | c c c c c c c c c c c @{}}
\toprule
$\Td/H_{i}$ &
$\{e\}$ & $\Z_{2}^{t}$ & $\Z_{2}^{d}$ & $\Z_{3}$ & $V_{4}$ & $V_{4}^{n}$
& $\Z_{4}$ & $S_{3}$ & $D_{4}$ & $A_{4}$ & $S_{4}$ \\
\midrule
$\Td/\{e\}$            & $24$ & $0$ & $0$ & $0$ & $0$ & $0$ & $0$ & $0$ & $0$ & $0$ & $0$ \\
$\Td/\Z_{2}^{t}$       & $12$ & $2$ & $0$ & $0$ & $0$ & $0$ & $0$ & $0$ & $0$ & $0$ & $0$ \\
$\Td/\Z_{2}^{d}$       & $12$ & $0$ & $4$ & $0$ & $0$ & $0$ & $0$ & $0$ & $0$ & $0$ & $0$ \\
$\Td/\Z_{3}$           & $8$  & $0$ & $0$ & $2$ & $0$ & $0$ & $0$ & $0$ & $0$ & $0$ & $0$ \\
$\Td/V_{4}$            & $6$  & $2$ & $2$ & $0$ & $2$ & $0$ & $0$ & $0$ & $0$ & $0$ & $0$ \\
$\Td/V_{4}^{n}$        & $6$  & $0$ & $6$ & $0$ & $0$ & $6$ & $0$ & $0$ & $0$ & $0$ & $0$ \\
$\Td/\Z_{4}$           & $6$  & $0$ & $2$ & $0$ & $0$ & $0$ & $2$ & $0$ & $0$ & $0$ & $0$ \\
$\Td/S_{3}$            & $4$  & $2$ & $0$ & $1$ & $0$ & $0$ & $0$ & $1$ & $0$ & $0$ & $0$ \\
$\Td/D_{4}$            & $3$  & $1$ & $3$ & $0$ & $1$ & $3$ & $1$ & $0$ & $1$ & $0$ & $0$ \\
$\Td/A_{4}$            & $2$  & $0$ & $2$ & $2$ & $0$ & $2$ & $0$ & $0$ & $0$ & $2$ & $0$ \\
$\Td/S_{4}$            & $1$  & $1$ & $1$ & $1$ & $1$ & $1$ & $1$ & $1$ & $1$ & $1$ & $1$ \\
\bottomrule
\end{tabular}
\end{table}

\subsection{Coset decomposition -- algebra matches geometry}

For every known dissection the Burnside ring element is a single transitive
$\Td$-set:

\begin{table}[H]
\centering\small
\caption{Each dissection is a single transitive $\Td$-set in $\Burn(\Td)$.}
\label{tab:coset-match}
\begin{tabular}{@{}c c c c @{}}
\toprule
$S_{n}$ & Coset space & Stabiliser & $\Burn(\Td)$-element \\
\midrule
$S_{1}$  & $\Td/S_{4}$        & $S_{4}$        & $[\Td/S_{4}]$        \\
$S_{2}$  & $\Td/A_{4}$        & $A_{4}$        & $[\Td/A_{4}]$        \\
$S_{3}$  & $\Td/D_{4}$        & $D_{4}$        & $[\Td/D_{4}]$        \\
$S_{4}$  & $\Td/S_{3}$        & $S_{3}$        & $[\Td/S_{3}]$        \\
$S_{6}$  & $\Td/V_{4}$        & $V_{4}$        & $[\Td/V_{4}]$        \\
$S_{8}$  & $\Td/\Z_{3}$       & $\Z_{3}$       & $[\Td/\Z_{3}]$       \\
$S_{12}$ & $\Td/\Z_{2}^{t}$   & $\Z_{2}^{t}$   & $[\Td/\Z_{2}^{t}]$   \\
$S_{24}$ & $\Td/\{e\}$        & $\{e\}$        & $[\Td/\{e\}]$        \\
\bottomrule
\end{tabular}
\end{table}

\subsection{Geometric vs.\ algebraic stabilisers}\label{ssec:geom-vs-alg}

We now make precise the relation between the two notions of
symmetry attached to the atlas: the piece-type stabiliser of
\Cref{def:piece-type} (recorded in \cref{tab:td-action}, ``algebraic'')
and the partition-level stabiliser $G_{\mathrm{diss}}$ of
\Cref{def:gdiss} (recorded in \cref{tab:gdiss}, ``geometric'').

\begin{lemma}[Geometric vs.\ algebraic stabilisers]\label{lem:geom-vs-alg-stab}
For each $S_{n}$ in the atlas let $P_{0}$ be the canonical piece of
\cref{tab:atlas-reps}, let $H_{n} := \Stab_{\Td}(P_{0})$ be its
piece-type stabiliser (\Cref{def:piece-type}, $|H_{n}| = 24/n$), and let
$G_{\mathrm{diss}}(S_{n})$ be its geometric symmetry group
(\Cref{def:gdiss}).  Then:
\begin{enumerate}
\item $G_{\mathrm{diss}}(S_{n})$ is the subgroup of $\Td$ listed in
      \cref{tab:gdiss}, and the orbit--stabiliser identity
      $|G_{\mathrm{diss}}(S_{n})| = n \cdot |\Stab_{0}|$
      holds for every $n$, where
      $\Stab_{0} = G_{\mathrm{diss}}(S_{n}) \cap H_{n}$ is the stabiliser
      of $P_{0}$ inside the geometric group.
\item $G_{\mathrm{diss}}(S_{n}) = \Td$ (i.e.\ $S_{n}$ is $\Td$-invariant
      in the sense of \Cref{def:td-invariance}) for
      $n \in \{1,12,24\}$, and only for these~$n$.
\item For $n \in \{2,3,4,6,8\}$ one has
      $G_{\mathrm{diss}}(S_{n}) \subsetneq \Td$; the elements of
      $\Td \setminus G_{\mathrm{diss}}(S_{n})$ send $S_{n}$ to a
      \emph{distinct} but congruent partition of $T$.
\item In every case $S_{n}$ is $G_{\mathrm{diss}}$-transitive
      (\Cref{def:gdiss-transitive}), i.e.\ the $n$ pieces form a single
      orbit under $G_{\mathrm{diss}}(S_{n})$; this is the statement
      certified in \Cref{thm:transitivity}.
\end{enumerate}
\end{lemma}

\begin{proof}
The numerical data of \cref{tab:gdiss} are verified by direct
computation in the embedding~\eqref{eq:atlas-coords}: for each $S_{n}$
we enumerate the $24$ elements $g \in \Td$ and test whether the image
partition $g(S_{n})$ equals $S_{n}$ setwise; the set of such $g$ is
$G_{\mathrm{diss}}(S_{n})$.  Items~(2) and~(3) are read off the
table.  The orbit--stabiliser identity in~(1) is the usual one for the
transitive action of $G_{\mathrm{diss}}(S_{n})$ on the $n$ pieces of
$S_{n}$; transitivity within $G_{\mathrm{diss}}(S_{n})$ is verified case
by case by the atlas computation.  The public script
\path{scripts/independent_verifier.py} gives a separate check in the stated
Coxeter-pure $15$-family scope.  Item~(4) is a restatement of the definitions
once transitivity of the $G_{\mathrm{diss}}$-action is established.
\end{proof}

\begin{remark}[Summary]\label{rem:sym-summary}
For $n \in \{1,12,24\}$ all four notions of \Cref{sec:prelim}
coincide; for $n \in \{2,3,4,6,8\}$ only $G_{\mathrm{diss}}$-transitivity
(\Cref{def:gdiss-transitive}) holds, and
$G_{\mathrm{diss}}(S_{n}) \subsetneq \Td$ is strictly smaller than
$\Td$.  In particular the Lagrange obstruction $n \mid 24$
(\Cref{cor:forward}) uses only $G_{\mathrm{diss}}$-transitivity.
\end{remark}

\begin{table}[H]
\centering\small
\caption{Geometric vs.\ algebraic stabilisers of the eight dissections.}
\label{tab:gdiss}
\begin{tabular}{@{} c c c c c c @{}}
\toprule
$S_{n}$ & $G_{\mathrm{diss}}$ & $|G_{\mathrm{diss}}|$
        & Transitive? & $|\Stab_{0}|$ \\
\midrule
$S_{1}$  & $S_{4}$ & $24$ & \checkmark & $24$ \\
$S_{2}$  & $V_{4}$ & $4$  & \checkmark & $2$  \\
$S_{3}$  & $S_{3}$ & $6$  & \checkmark & $2$  \\
$S_{4}$  & $D_{4}$ & $8$  & \checkmark & $2$  \\
$S_{6}$  & $S_{3}$ & $6$  & \checkmark & $1$  \\
$S_{8}$  & $D_{4}$ & $8$  & \checkmark & $1$  \\
$S_{12}$ & $S_{4}$ & $24$ & \checkmark & $2$  \\
$S_{24}$ & $S_{4}$ & $24$ & \checkmark & $1$  \\
\bottomrule
\end{tabular}
\end{table}

%%%%%%%%%%%%%%%%%%%%%%%%%%%%%%%%%%%%%%%%%%%%%%%%%%%%%%%%%%%%%%%%%%%%%%%%
%%%%%%%%%%%%%%%%%%%%%%%%%%%%%%%%%%%%%%%%%%%%%%%%%%%%%%%%%%%%%%%%%%%%%%%%
\section{IFS dissection grammar}\label{sec:ifs}

Every halving refinement in the poset of \Cref{sec:poset,sec:refinement} can
be encoded as an iterated function system (IFS): a pair of affine
contractions whose images tile the parent piece.  This produces a formal
grammar for generating all eight dissections.

\subsection{Halving rules}

\begin{table}[H]
\centering\small
\caption{Halving rules and their chirality.  Rule R3 ($S_{4} \to S_{8}$) is
the unique rule with \emph{both} branches orientation-preserving.}
\label{tab:ifs-rules}
\begin{tabular}{@{} c l l c c @{}}
\toprule
Rule & Production & Geometric name & $\det(f_{1})$ & $\det(f_{2})$ \\
\midrule
R1 & $S_{1} \to f_{1}(S_{2}) \cup f_{2}(S_{2})$          & Median bisection            & $+1$ & $-1$ \\
R2 & $S_{2} \to g_{1}(S_{4}) \cup g_{2}(S_{4})$          & Orthogonal median           & $+1$ & $-1$ \\
R3 & $S_{4} \to h_{1}(S_{8}) \cup h_{2}(S_{8})$          & Longest-edge bisection      & $+1$ & $+1$ \\
R4 & $S_{3} \to p_{1}(S_{6}) \cup p_{2}(S_{6})$          & Axial median bisection      & $+1$ & $-1$ \\
R5 & $S_{12} \to q_{1}(S_{24}) \cup q_{2}(S_{24})$       & Edge bisection              & $+1$ & $-1$ \\
\bottomrule
\end{tabular}
\end{table}

\subsection{Heterogeneous bridge (R6)}

The heterogeneous bridge (\cref{thm:s2-bridge}) corresponds to the rule
$R6 \colon S_{2} \to \alpha(S_{3}) \cup \beta(S_{6})$ with
$\det(\alpha)\det(\beta) = -1$.

\subsection{Formal alphabet and grammar}

The alphabet of the grammar is
$\Sigma = \{S_{1}, S_{2}, S_{3}, S_{4}, S_{6}, S_{8}, S_{12}, S_{24}\}$.
Homogeneous productions are R1--R5; the unique heterogeneous production is
R6.  Derived multi-level compositions include $R1 \circ R2$
($S_{1} \to 4 \times S_{4}$, quadrisection) and $R1 \circ R2 \circ R3$
($S_{1} \to 8 \times S_{8}$, $8T$-LE octasection).

\subsection{Composition \texorpdfstring{$S_{1} \to 8 \cdot S_{8}$}{S1 -> 8 S8}
and PPS23 similarity classes}

Iterated longest-edge bisection of $T$ produces a finite sequence of
similarity classes~\cite{PPS2023}; our grammar reproduces and
extends their enumeration:

\begin{table}[H]
\centering\small
\caption{Similarity class count under iterated longest-edge bisection.}
\label{tab:pps23}
\begin{tabular}{@{} r r r @{}}
\toprule
Depth & Pieces & Similarity classes \\
\midrule
$1$  & $2$   & $1$ \\
$2$  & $4$   & $2$ \\
$3$  & $8$   & $4$ \\
$4$  & $16$  & $8$ \\
$5$  & $32$  & $9$ \\
$6$  & $64$  & $12$ \\
$7$  & $128$ & $16$ \\
$8$  & $256$ & $17$ \\
$9$  & $512$ & $\mathbf{18}$ \\
\bottomrule
\end{tabular}
\end{table}

The class count stabilises at $18$ by depth $9$, consistent with a finite
attractor in similarity space.  At depth $9$ the four dominant classes are
the regular tetrahedron ($\times 40$), an isosceles variant ($\times 40$),
the $S_{8}$-type ($\times 108$) and the $S_{3}$-type ($\times 40$).

%%%%%%%%%%%%%%%%%%%%%%%%%%%%%%%%%%%%%%%%%%%%%%%%%%%%%%%%%%%%%%%%%%%%%%%%
\section{An obstruction stack for $\Td$-invariant dissections
(screening)}\label{sec:nontrans}

\noindent\textbf{Status of this section (honesty block).}  This section
does \emph{not} resolve the non-$G_{\mathrm{diss}}$-transitive case of
\cref{conj:divisor}.  It builds an \emph{evidence stack} that, under the
standing hypothesis of $\Td$-invariance (\Cref{def:td-invariance}) and a
$\Td$-maximality argument flagged as \emph{plausible but not fully
rigorous} (\Cref{ssec:td-max,rem:td-max-caveat}), filters every tested
orbit decomposition for $n \leq 16$, $n \nmid 24$.  Of the nine such
values of $n$ four pass through the crystallographic restriction alone
(\emph{unconditional} screening in the $\Td$-invariant category) and five
additionally invoke the $\Td$-maximality step, tagged~$(\dagger)$ in
\cref{tab:orbit-filter} and recorded at level \textsf{L1+} in
\cref{tab:argument-chain}.  A theorem-level resolution of the
non-transitive problem, without the $\Td$-maximality caveat, is the
content of the Stabiliser Transport Lemma
(\Cref{lem:stab-transport}, \Cref{sec:programme}) and is not achieved
here.

\medskip

The remaining question for \cref{conj:divisor} is whether a congruent
$n$-dissection of $T$ with $n \nmid 24$ could exist without being
$G_{\mathrm{diss}}$-transitive (\Cref{def:gdiss-transitive}).  We develop a
systematic obstruction framework combining the subgroup lattice, the
crystallographic restriction theorem, and a $\Td$-maximality argument; the
resulting stack filters all tested decompositions in the $\Td$-invariant
category (\Cref{def:td-invariance}).  It is important to stress that
$\Td$-invariance is strictly stronger than mere non-transitivity under
$G_{\mathrm{diss}}$: a hypothetical congruent dissection might have
$G_{\mathrm{diss}}(D) \subsetneq \Td$ and therefore escape the $\Td$-orbit
analysis of this section entirely.

\subsection{Orbit decomposition}\label{ssec:orbit-decomp}

Suppose a congruent $n$-dissection $T = \bigsqcup_{i=1}^{n} P_{i}$ is
\emph{$\Td$-invariant} (\Cref{def:td-invariance}), i.e.\ every
$\sigma \in \Td$ permutes the pieces
(equivalently, $\sigma(\{P_1,\ldots,P_n\}) = \{P_1,\ldots,P_n\}$ for all
$\sigma$).  All eight known dissections enjoy this property for
$n \in \{1,12,24\}$ (\cref{lem:geom-vs-alg-stab}); for
$n \in \{2,3,4,6,8\}$ only the smaller group
$G_{\mathrm{diss}}(S_{n}) \subsetneq \Td$ preserves the partition set-wise
and the analysis below strictly speaking applies only after replacing $\Td$
by $G_{\mathrm{diss}}(S_{n})$.  The present section enforces $\Td$-invariance
as a standing hypothesis so as to have the full group $\Td$ available for
the orbit decomposition.
Under this assumption, the $\Td$-action decomposes into orbits
$\calO_{1}, \ldots, \calO_{k}$ with
$|\calO_{j}| = [\Td : H_{j}]$ for stabiliser $H_{j} \leq \Td$.  Since each
orbit size divides $|\Td| = 24$:
\begin{equation}\label{eq:orbit-decomp}
  n = \sum_{j=1}^{k}\, [\Td : H_{j}],
  \qquad
  [\Td : H_{j}] \in \{1,2,3,4,6,8,12,24\}.
\end{equation}
For $n \nmid 24$ necessarily $k \geq 2$.

\begin{lemma}[Size-$1$ orbit]\label{lem:size-1}
At most one orbit has size $1$.
\end{lemma}

\begin{proof}
A $\Td$-invariant piece must contain the centroid $G_{T}$, the unique
$\Td$-fixed point of $T$.  Two distinct convex pieces containing $G_{T}$
would overlap.
\end{proof}

\subsection{Cross-orbit symmetry obstruction}

Let $P_{1} \in \calO_{1}$ and $P_{2} \in \calO_{2}$ be pieces in distinct
orbits of a congruent $\Td$-invariant dissection.  Since $P_{1} \cong P_{2}$,
a fixed congruence identifies them with a common abstract reference tile
$P$ whose intrinsic symmetry group is $K := \Sym(P) \leq O(3)$.  The
stabilisers
\[
  H_{1} = \Stab_{\Td}(P_{1}), \qquad H_{2} = \Stab_{\Td}(P_{2})
\]
act on $P$ and hence embed (as subgroups, up to conjugation in $O(3)$) into~$K$.

\begin{lemma}[Stabiliser pair screening]\label{lem:stab-embed}
Let $H_{1}, H_{2} \leq \Td$ be two point-stabiliser conjugacy classes from
the $30$ subgroups of $\Td$ (in $11$ conjugacy classes).  Let
$K(H_{1}, H_{2}) \leq O(3)$ denote the smallest finite subgroup of
$O(3)$, up to conjugation, that contains a pair of non-conjugate copies of
$H_{1}$ and $H_{2}$ (whenever such a pair exists).  Then, using the
classical classification of finite subgroups of $O(3)$ and the
crystallographic restriction theorem~\cite{Hiller1985}, the following
non-conjugate stabiliser pairs force $K$ to contain a full cubic group
$\geq S_{4}$:
\begin{center}
\begin{tabular}{@{}ll@{}}
\toprule
Non-conjugate pair $(H_{1}, H_{2})$ & Minimal containing $K \leq O(3)$ \\
\midrule
$(D_{4}, \Z_{3})$      & $\supseteq S_{4}$ \quad (mixed $4$-fold \& $3$-fold) \\
$(\Z_{4}, \Z_{3})$      & $\supseteq S_{4}$ \quad (same mechanism) \\
$(S_{3}, V_{4})$       & $\supseteq S_{4}$ \quad ($3$-fold + orthogonal $2$-folds) \\
$(A_{4}, D_{4})$       & $\supseteq S_{4}$ \quad (automatic) \\
$(A_{4}, S_{3})$       & $\supseteq S_{4}$ \quad (automatic) \\
$(D_{4}, S_{3})$       & $\supseteq S_{4}$ \quad (automatic) \\
\bottomrule
\end{tabular}
\end{center}
In particular, for each such pair, any abstract tile $P$ admitting both
$H_{1}$ and $H_{2}$ as non-conjugate subgroups of $\Sym(P)$ satisfies
$|\Sym(P)| \geq 24$ and $\Sym(P)$ contains a copy of $S_{4}$ acting on a
cube.
\end{lemma}

\begin{proof}[Proof sketch]
Each listed pair either (i) combines a $4$-fold and a $3$-fold rotation,
or (ii) combines a $3$-fold rotation with a $V_{4}$ subgroup of orthogonal
$2$-folds, or (iii) already contains $A_{4}$.  In case~(i) the
crystallographic restriction and the classification of finite rotation
groups of $\R^{3}$ force the generated subgroup to be the rotation group
of the cube (order $24$) or its full symmetry group (order $48$); in
case~(ii) the three mutually orthogonal $2$-folds and the $3$-fold
permuting them generate the same rotation subgroup; case~(iii) is
immediate since $A_{4} \leq S_{4}$.  A finite verification on all unordered
pairs of conjugacy classes in $\Td$ (performed computationally; see
\path{scripts/independent_verifier.py} in the public companion repository)
confirms no other non-conjugate pair realises a smaller containing group.
\end{proof}

\begin{remark}[Why this is a screening lemma, not an impossibility theorem]
\label{rem:stab-screen}
\cref{lem:stab-embed} is used as a \emph{finite screening device} on orbit
decompositions of $n \nmid 24$: whenever a putative decomposition assigns
stabilisers in a pair of the table above to distinct orbits, the abstract
tile must admit a cubic-or-larger symmetry group, which is incompatible
with the finite list of tile shapes compatible with a tetrahedral
dissection (verified case-by-case in the obstruction stack,
\cref{tab:orbit-filter}).  It is \emph{not} a standalone impossibility
theorem; the final kill still requires either $\Td$-maximality
(\cref{ssec:td-max}, with the caveat of \cref{rem:td-max-caveat}) or the
Stabiliser Transport Lemma (\cref{lem:stab-transport}).
\end{remark}

\subsection{The \texorpdfstring{$\Td$}{Td}-maximality argument}\label{ssec:td-max}

\noindent\emph{Status.}  The argument of this subsection is \textbf{plausible
but not fully rigorous}: it relies on the implicit assumption that the
cross-orbit congruence $P_{1} \to P_{2}$ extends to a global isometry of
$T$, which is not forced by the dissection hypothesis alone.  See
\cref{rem:td-max-caveat} immediately below and the Stabiliser Transport
Lemma (\cref{lem:stab-transport}, \cref{sec:programme}) for the correct
group-theoretic completion.  Decompositions marked $(\dagger)$ in
\cref{tab:orbit-filter} are precisely those whose filtering depends on
this conditional argument.

\medskip

For orbit decompositions surviving the crystallographic filter -- those with
\emph{same-type} stabilisers (e.g.\ $n = 9 = 3 + 3 + 3$ with all $D_{4}$) --
the cross-orbit congruence map $P_{1} \to P_{2}$ must be an isometry of
$\R^{3}$ that is \emph{not} in $\Td$: $\Td$ permutes pieces within orbits,
not between them.  Since $\Td = \Isom(T)$ is the full isometry group of $T$,
any such map cannot preserve $T$ globally.

\begin{remark}[Caveat on \Cref{ssec:td-max}]\label{rem:td-max-caveat}
A congruence $\sigma \colon P_{i} \to P_{j}$ only needs to map one piece to
another; it does not need to preserve $T$ globally.  Hence
$\sigma \notin \Td$ does not \emph{itself} produce a contradiction: the
argument of \Cref{ssec:td-max} is plausible but not fully rigorous, and the
Stabiliser Transport Lemma of \Cref{sec:programme} (\cref{lem:stab-transport})
provides the correct group-theoretic completion.
\end{remark}

\subsection{Results of the obstruction stack}

Systematic enumeration of all integer partitions of $n$ into valid orbit
sizes (\eqref{eq:orbit-decomp}), followed by stabiliser-type analysis, gives:

\begin{table}[H]
\centering\small
\caption{Non-transitive orbit decompositions filtered for $n \leq 16$,
$n \nmid 24$.  A dagger $(\dagger)$ indicates decompositions that require
the $\Td$-maximality argument of \Cref{ssec:td-max}, plausible but not yet
rigorous.}
\label{tab:orbit-filter}
\begin{tabular}{@{} c c c c l @{}}
\toprule
$n$ & Total decomps. & After filter & All filtered? & Primary mechanism \\
\midrule
$5$   & $6$   & $1$  & yes             & $K \geq S_{4}$ \\
$7$   & $12$  & $3$  & yes             & $K \geq S_{4}$ \\
$9$   & $22$  & $6$  & yes$^{\dagger}$ & $K \geq S_{4}$ + same-type + $\Td$-max \\
$10$  & $30$  & $8$  & yes$^{\dagger}$ & $K \geq S_{4}$ + same-type + $\Td$-max \\
$11$  & $36$  & $10$ & yes             & $K \geq S_{4}$ \\
$13$  & $58$  & $17$ & yes             & $K \geq S_{4}$ \\
$14$  & $76$  & $21$ & yes$^{\dagger}$ & $K \geq S_{4}$ + same-type + $\Td$-max \\
$15$  & $90$  & $26$ & yes$^{\dagger}$ & $K \geq S_{4}$ + same-type + $\Td$-max \\
$16$  & $115$ & $33$ & yes$^{\dagger}$ & $K \geq S_{4}$ + same-type + $\Td$-max \\
\bottomrule
\end{tabular}
\end{table}

\noindent
Four of the nine cases are rigorously filtered by the crystallographic
restriction alone (marked without dagger in \cref{tab:orbit-filter});
the remaining five additionally require the
$\Td$-maximality argument (\Cref{ssec:td-max}), which is plausible but not
yet fully rigorous -- see \cref{rem:td-max-caveat}.  We therefore do
\emph{not} claim theorem-level impossibility of non-$G_{\mathrm{diss}}$-transitive
dissections for $n \leq 16$: what the table establishes is that, within
the $\Td$-invariant category, every tested orbit decomposition is ruled
out \emph{up to} the $\Td$-maximality caveat.  This is recorded at level
\textsf{L1+} in \cref{tab:argument-chain} and motivates (but does not
prove) \cref{conj:transitivity}.

\subsection{Argument chain summary}

\begin{table}[H]
\centering\small
\caption{Evidence-level stack of the screening argument against
non-$G_{\mathrm{diss}}$-transitive dissections, under the standing
hypothesis of $\Td$-invariance.}
\label{tab:argument-chain}
\begin{tabular}{@{} c l c @{}}
\toprule
Step & Statement & Status \\
\midrule
1 & Non-transitive $\Rightarrow$ $\geq 2$ orbits, sizes divide $24$       & \textbf{Theorem} \\
2 & At most one $\Td$-invariant piece (\cref{lem:size-1})                 & \textbf{Theorem} \\
3 & Stabiliser action on vertices forces $|V(P) \cap V(T)| \in \{0, 4\}$ & Proposition (convex) \\
4 & Mixed stabiliser types $\Rightarrow$ $K \geq S_{4}$
       (crystallographic restriction)                                      & \textbf{Theorem} \\
5 & Same-type orbits require non-$\Td$ isometry                          & Plausible (L1+) \\
6 & $S_{4}$-symmetric bodies cannot tile $T$ for $n \nmid 24$            & L1 evidence \\
\bottomrule
\end{tabular}
\end{table}

%%%%%%%%%%%%%%%%%%%%%%%%%%%%%%%%%%%%%%%%%%%%%%%%%%%%%%%%%%%%%%%%%%%%%%%%
\section{The rigidity and configurational programme}\label{sec:programme}

We now describe the programme that attacks \cref{conj:transitivity} and
contains the two principal conditional configurational results of the paper:
\cref{thm:n5} (the $n = 5$ BRH conditional obstruction) and
\cref{thm:n16-aof-conditional} ($n = 16$ all-one-face conditional
impossibility).

Throughout this section we adopt two standing hypotheses, consistent with
all prior literature on congruent simplicial tilings: (i) every piece is
convex, and (ii) the dissection is \emph{face-to-face} (every interior
$2$-face of the piece complex is shared by exactly two pieces).

\subsection{Dihedral rigidity programme}\label{ssec:rigidity}

A direct attack via dihedral-angle analysis exploits the irrationality of
$\alpha/\pi$ (Conway--Jones~\cite{ConwayJones1976}).  The intended chain was:
$\alpha/\pi \notin \Q \Rightarrow$ interior edges carry zero
$\alpha$-coefficient $\Rightarrow$ Dehn localises to the original $T$-edges
$\Rightarrow$ every piece is an $\alpha$-carrier $\Rightarrow$ every piece
touches $\geq 2$ of $T$'s $4$ faces $\Rightarrow$ Face-Normal Rigidity forces
the gluing isometry $\phi \in \Td$ $\Rightarrow$ transitivity
$\Rightarrow$ $n \mid 24$.

\begin{theorem}[Face-Normal Rigidity]\label{thm:fnr}
Let $Q \in O(3)$ be an orthogonal linear map.  If $Q$ sends at least two
of the four face normals of~$T$ (as unit vectors based at the origin) to
face normals of~$T$, then $Q$ is the linear part of some element
of~$\Td$.  Equivalently, for every isometry
$\phi(x) = Qx + t$ of $\R^{3}$ whose linear part~$Q$ sends two or more
face normals of~$T$ to face normals of~$T$, one has $Q \in \Td$ (regarded
as a subgroup of $O(3)$ in the standard way).  Verified exhaustively on
all $144$ ordered pairs of face normals.
\end{theorem}

However the chain above is broken by two identified gaps:

\begin{description}
\item[GAP A (Dehn localisation).]
  The chain assumes all piece dihedral angles have the form
  $\theta = m\alpha + q\pi$ with $m \in \Z$.  This is unjustified for
  arbitrary cuts: only original $T$-edges carry $\alpha$-related angles a
  priori, and global cancellation $\sum m_{i} = 0$ at an interior edge
  does \emph{not} force each $m_{i}$ to vanish.  The na\"ive
  ``$m \geq 2 \Rightarrow \theta > \alpha$'' dichotomy also fails, e.g.\
  $\theta = 2\alpha - \pi/2 \approx 51^{\circ} < \alpha$.
\item[GAP B (Intrinsic boundary).]
  A face of a piece lying on $\partial T$ is not intrinsically
  distinguishable from an interior face: the gluing isometry
  $\phi \colon P_{1} \to P_{2}$ could map a boundary face of $P_{1}$ to an
  interior face of $P_{2}$.  Without intrinsic distinguishability the
  Face-Normal Rigidity theorem (\cref{thm:fnr}, which \emph{is} valid as a
  finite test) cannot be invoked to force $\phi \in \Td$.
\end{description}

\cref{thm:fnr} is therefore a correct rigidity tool, but its deployment
requires first discharging GAPs A and B.  The remainder of \Cref{sec:programme}
addresses these gaps.

\subsection{Reduction engine and Stabiliser Transport}\label{ssec:reduction-engine}

Following the GAP A/B diagnosis we pivot to a \emph{reduction strategy}:
defer the final kill and restrict the shape of any putative counterexample.

\begin{lemma}[Stabiliser Transport -- \textsf{L2} theorem]\label{lem:stab-transport}
Let $T = P_{1} \cup \cdots \cup P_{n}$ be a congruent dissection with
congruence group $G = \langle \sigma_{ij} \rangle \leq \Isom(\R^{3})$.
Then:
\begin{enumerate}
  \item[(a)] stabilisers of congruent pieces are conjugate in $\Isom(\R^{3})$;
  \item[(b)] if $n \nmid 24$ then $G \not\leq \Td$; consequently $G$ contains
             a non-tetrahedral congruence and the $\Td$-action on the pieces
             has $\geq 2$ orbits.
\end{enumerate}
\end{lemma}

\cref{lem:stab-transport}(a) is standard; (b) follows from the fact that
$\Td$ has $30$ subgroups in $11$ conjugacy classes with index set
$\{24/|H|\} = \mathrm{div}(24)$.

\begin{definition}[Boundary-Normal set]\label{def:bn}
For each face $f$ of the abstract tile $P$, the \emph{boundary-normal set}
\[
  \BN(f) \coloneqq \bigl\{\, i \in \{0,1,2,3\}
          : \exists\ \text{placement } \phi \text{ with }
          \phi(f) \subset F_{i}(T) \,\bigr\}
\]
records which $T$-faces $F_{0},\ldots,F_{3}$ the abstract face $f$ can land
on.
\end{definition}

\begin{conjecture}[Boundary-Normal Orbits, L1]\label{conj:bn-orbits}
In a minimal counterexample to \cref{conj:transitivity}, the boundary faces
of the abstract tile $P$ fall into at most two $\Sym(P)$-orbits under the
$4$-normal dictionary $\BN(\,\cdot\,)$.
\end{conjecture}

The face-orbit analysis for all $11$ subgroup conjugacy classes of $\Td$
partitions $\{0,1,2,3\}$ into progressively finer schemes: e.g.\ $\Z_{3}$
acts as $\{0\}\,|\,\{1,2,3\}$, $V_{4}$ as $\{0,1\}\,|\,\{2,3\}$, which
constrains the boundary structure of any counterexample.

\subsection{Spherical Link Program and the Conditional Separation Lemma}
\label{ssec:spherical-link}

The local geometry at each $T$-vertex $V$ is captured by the \emph{vertex
link}, an equilateral spherical triangle $\Delta_{V}$.  Its angles are equal
to the dihedral $\alpha$ of $T$, and its sides equal the face-angle $\pi/3$:
both are exactly computable.

\begin{proposition}[Vertex-link geometry, L2]\label{prop:vertex-link}
The vertex link $\Delta_{V}$ at every vertex of $T$ is equilateral with
spherical angle $\alpha = \arccos(1/3)$ and side length $\pi/3$.
\end{proposition}

\begin{proof}
Apply the spherical law of cosines to the three faces meeting at $V$:
$\cos a = \cos\alpha/(1 - \cos\alpha) = 1/2$, whence $a = \pi/3$.  Equality of
the three angles is immediate from the symmetry $\Sym_{V}(T)$.
\end{proof}

\begin{proposition}[Conditional Separation Lemma, L1+]\label{prop:csl}
Suppose \emph{Hypothesis H} holds: every interior angle of every link-tile
at every $T$-vertex belongs to $\Q\alpha + \Q\pi$.
Then under the Gauss--Bonnet constraint on $s$-gon tiles $Q$ of
$\Delta_{V}$, the identity separates into
\[
  \sum_{\ell} a_{\ell} = \frac{3}{k},
  \qquad
  \sum_{\ell} b_{\ell} = s - 2 - \frac{1}{k},
  \qquad a_{\ell}, b_{\ell} \in \Q,
\]
where $k$ is the number of pieces meeting at $V$.  In a uniform convex model
this forces $k \in \{1, 3, 6, 9, 12, \ldots\}$, eliminating $k \in \{2, 4, 5\}$.
\end{proposition}

Three open gaps prevent \cref{prop:csl} from attaining theorem status:
\textbf{GAP C} (Hypothesis H), \textbf{GAP D} (mixed vertex types may
produce non-congruent link-tiles), \textbf{GAP E} (non-convex pieces).
Discharge of GAP C is the key conditional step.

%%%%%%%%%%%%%%%%%%%%%%%%%%%%%%%%%%%%%%%%%%%%%%%%%%%%%%%%%%%%%%%%%%%%%%%%
\subsection{Configurational closure at \texorpdfstring{$n = 16$}{n=16}:
the all-one-face residual}\label{sec:prog-n16}

\cref{conj:bn-orbits} targets GAP B by classifying boundary-face orbits
abstractly.  The same gap can be attacked \emph{configurationally}:
enumerate, for fixed piece count $n$, the candidate ``abstract role
distributions'' of pieces on $\partial T$ and eliminate them one by one.

For $n = 16$ the most constrained distribution places one boundary face of
every piece on the \emph{same} single $T$-face; we call this the
\emph{all-one-face} residual.

\paragraph{Setup.}
Embed $T$ via the rational coordinates
\begin{equation}\label{eq:16-embedding}
  v_{0} = (1,1,1),\quad
  v_{1} = (1,-1,-1),\quad
  v_{2} = (-1,1,-1),\quad
  v_{3} = (-1,-1,1),
\end{equation}
so $|v_{i} v_{j}|^{2} = 8$ and $V_{T} = 8/3$.  The medial triangle of each
$T$-face decomposes that face into four sub-triangles (\emph{slots}): three
corner slots (at the three $T$-vertices) and one central slot (the medial
triangle itself).  The conditional theorem uses the following two explicit
finite hypotheses to cut the search space:

\begin{description}
\item[$(\Hroles)$] (Role-type assignment).
  The medial decomposition is the only abstract role-type assignment
  compatible with face-to-face adjacency: each boundary footprint is either
  a corner slot or the central slot.
\item[$(\Hslot)$] (Slot separation).
  Every piece's boundary footprint is exactly one of the four slots above;
  no piece's footprint mixes slot types.
\end{description}

Both $(\Hroles)$ and $(\Hslot)$ are unconditional in the
\emph{Coxeter-pure} instance (where every piece is a union of
$A_{3}$-orthoschemes) but residual in the general convex instance.

\begin{theorem}[$n=16$ all-one-face: Coxeter-pure instance]\label{thm:n16-aof-coxeter}
There is no face-to-face Coxeter-pure congruent dissection of $T$ into
$16$ mutually congruent pieces such that every piece has exactly one
boundary face on a single common $T$-face.
\end{theorem}

\begin{proof}
In the Coxeter-pure instance both $(\Hroles)$ and $(\Hslot)$ are automatic
(a union of $A_{3}$-orthoschemes has a unique role assignment and a
unique boundary-slot type per piece; see \Cref{sec:prog-n16}).  The
conclusion therefore follows from the conditional
\cref{thm:n16-aof-conditional} below.
\end{proof}

\begin{conditionaltheorem}[$n=16$ all-one-face: general convex instance]\label{thm:n16-aof-conditional}
Assume $(\Hroles) + (\Hslot)$.  Let $\{P_{1},\ldots,P_{16}\}$ be a
face-to-face congruent dissection of $T$ into $16$ mutually congruent
convex polytopes such that every $P_{i}$ has exactly one boundary face on a
single common $T$-face.  Then no such dissection exists.
\end{conditionaltheorem}

The proof assembles three rigidity lemmas, each independently certified by
exact rational arithmetic (Python \texttt{Fractions}, zero failed tests
across the chain).

\begin{lemma}[Tetrahedrality reduction, C24x]\label{lem:c24x}
Every $P_{i}$ is a tetrahedron with triangular medial base on $\partial T$
and a single interior apex.
\end{lemma}

\begin{proof}
Two parts.

\emph{(I) Scaled-top family.}
No \emph{scaled-top} piece -- right prism ($t = 1$), right frustum
($t \in (0,1)$), or right pyramid with apex above the base centroid ($t = 0$)
-- fits.  The volume constraint $V_{\mathrm{piece}} = V_{T}/16 = 1/6$ with base
area $\sqrt 3/2$ forces height $h_{\mathrm{piece}} = h_{T}/4$ (prism) or the
frustum-corrected analogue.
\begin{itemize}
  \item \emph{Layer A (prism).}  At $h_{\mathrm{prism}} = \sqrt 3/9$ the
        $T$-cross-section has squared inradius $(121/144)(2/3)$, while the
        medial triangle has squared circumradius $2/3$; the ratio
        $144/121 > 1$, so at least one medial vertex sticks out of $T$.
  \item \emph{Layer B (frustum).}  For any base point $Q$ the difference
        between the corner-frustum top and the central-frustum top equals
        $(1 - t)\,(c_{\mathrm{corner}} - c_{\mathrm{central}})$.  Since
        $|c_{\mathrm{corner}} - c_{\mathrm{central}}|^{2} = 2/3 \neq 0$
        exactly, face-to-face matching on the shared medial edge fails for
        every $t \in [0, 1)$.
\end{itemize}

\emph{(II) Multi-apex residual.}
Every convex piece with $k \geq 2$ interior vertices is excluded by two
independent rigorous arguments.  Let $F_{1}$ be the lateral face of a
corner-$v_{j}$ piece carrying the medial edge $(m_{jk}, m_{jl})$, and let
$q$ be the number of interior vertices of the piece on $F_{1}$.
\begin{itemize}
  \item \textbf{(ARG-1) $F_{4}$ parity-pairing} ($k = 2, q = 2$, Type-$2$
        $5$-vertex polytope).  The apex-interior face
        $F_{4} = (v_{j}, A, B)$ contains $v_{j}$ and both interior vertices
        but no midpoint.  At vertex $v_{j}$ exactly three corner-$v_{j}$
        pieces meet, each carrying one such $F_{4}$.  Face-to-face matching
        forces these three faces to pair \emph{among themselves}, but a
        fixed-point-free involution on a $3$-element set is impossible.
  \item \textbf{(ARG-2) Three-plane collapse} (all $k \geq 2$, $q \geq 2$).
        Congruence on edge multisets forces every interior vertex $X$ on
        $F_{1}$ of a corner-$v_{1}$ piece to satisfy
        $|v_{1} - X|^{2} = |m_{kl} - X|^{2}$, i.e.\ $X$ lies on the
        perpendicular bisector plane of $(v_{1}, m_{kl})$.  Face-to-face
        propagation across the three $T$-faces incident to $v_{1}$ yields
        three such planes:
        \[
          2x - y - z = 1, \qquad x - y - 2z = 1, \qquad x - 2y - z = 1.
        \]
        The coefficient matrix is non-singular with determinant $-4$; the
        unique intersection point is
        \[
          P_{*} = \tfrac{1}{4}(1, -1, -1),
        \]
        a rational interior point of $T$ with barycentric coordinates
        $(3, 7, 3, 3)/16$.  Every interior vertex on $F_{1}$ collapses to
        $P_{*}$, contradicting $q \geq 2$ distinct vertices.
  \item The case $q = 1$ (triangular $F_{1}$) reduces to \cref{lem:c24w}
        below via face-by-face apex collapse along the three medial edges.
\end{itemize}
This completes the proof.
\end{proof}

\begin{lemma}[Apex collapse, C24w]\label{lem:c24w}
Conditional on \cref{lem:c24x}, every congruent piece has a single interior
apex at perpendicular height $h_{T}/4$ above its medial base, and four
rigidity layers force all $16$ apices to coincide at the centroid
$G_{T} = (0, 0, 0)$:
\begin{itemize}
  \item[$(L_{1})$] Volume forces apex height $= h_{T}/4$.
  \item[$(L_{2})$] Same-face medial-edge matching forces the apex of the
        corner-$v_{j}$ piece on face $F_{i}$ to equal the apex of the
        central piece on $F_{i}$.
  \item[$(L_{3})$] $T$-edge matching propagates the apex across faces,
        forcing all four central apices to coincide at a common point $Q$.
  \item[$(L_{4})$] $Q$ is at perpendicular distance $h_{T}/4$ from every
        face; the $3$-of-$4$ face-normal linear system has unique solution
        $Q = G_{T}$.
\end{itemize}
\end{lemma}

\begin{lemma}[Edge-multiset incongruence at $G_{T}$, C24v]\label{lem:c24v}
At $Q = G_{T}$, the corner piece has squared-edge multiset
$\{1, 1, 2, 2, 2, 3\}$ and the central piece has squared-edge multiset
$\{1, 1, 1, 2, 2, 2\}$.  These multisets differ; hence the corner and
central pieces are not congruent.
\end{lemma}

\begin{proof}[Proof of \cref{thm:n16-aof-conditional}]
Combine \cref{lem:c24x}, \cref{lem:c24w}, and \cref{lem:c24v} with the
medial-base reduction encoded in $(\Hroles) + (\Hslot)$: the first forces
tetrahedrality; the second collapses all apices to $G_{T}$; the third
produces a direct violation of piece congruence.
\end{proof}

\begin{remark}[Scope of \cref{thm:n16-aof-coxeter,thm:n16-aof-conditional}]
In the Coxeter-pure instance $(\Hroles) + (\Hslot)$ are automatic,
so \cref{thm:n16-aof-coxeter} is the \textsf{L3} (unconditional)
specialisation of the \textsf{L2} conditional statement
\cref{thm:n16-aof-conditional}.  In the
general convex instance $(\Hroles) + (\Hslot)$ remain hypotheses: an
unconditional proof of them in that setting is not supplied here.  Thus
\cref{thm:n16-aof-conditional} is at evidence level \textsf{L2}.
\end{remark}

\begin{remark}[Convexity hypothesis]
The convexity of the pieces is essential to \cref{lem:c24x}: without it,
multi-apex non-convex pieces are not reached by the perpendicular-bisector
argument (ARG-$2$).  Discharging convexity is natural future work; the
present theorem is, as stated, in the convex category.
\end{remark}

\begin{remark}[Why $n = 16$?]
$n = 16$ divides neither $24$ nor $48$.  Note that closing the full
$n = 16$ case \emph{would} advance \cref{conj:divisor}, since $16 \nmid 24$
means that the conjecture predicts no dissection should exist.
However, \cref{lem:c24x} closes only the \emph{all-one-face residual}
configuration at $n = 16$, not the entire case.  This particular
configuration was historically the most resistant residual of the
role-distribution table in the programme described in \Cref{ssec:reduction-engine}.
The same C24 machinery is structurally portable to the other
non-divisors $n \in \{20, 28, \ldots\}$ and gives a template for the
configurational attack on GAP~B.
\end{remark}

%%%%%%%%%%%%%%%%%%%%%%%%%%%%%%%%%%%%%%%%%%%%%%%%%%%%%%%%%%%%%%%%%%%%%%%%
\subsection{The \texorpdfstring{$n=5$}{n=5} case: a conditional boundary-role obstruction}
\label{sec:prog-n5}

The previous public reduction of arbitrary convex face-to-face five-piece
dissections to one accidental-even metric residual is not established.  Its
boundary-role step requires an additional hypothesis: congruent copies need
not expose the same facets of the common reference polytope.  We therefore
replace that reduction by the conditional result below.  It is a complete
obstruction under an explicit boundary-role homogeneity condition (BRH),
while the unrestricted case with different boundary roles remains open.

\begin{definition}[Boundary-role homogeneity (BRH)]\label{def:brh}
Let $P_1,\ldots,P_5$ be full-dimensional congruent convex polytopes with
union $T$ and pairwise disjoint interiors.  BRH holds if there are a
reference copy $P$, Euclidean isometries $\gamma_i$ with
$\gamma_i(P)=P_i$, and a fixed subset $B$ of facets of $P$ such that
$\gamma_i(B)$ is exactly the set of facets of $P_i$ contained in
$\partial T$.  The set $B$ is the common boundary-role set.  BRH requires
congruences preserving these designated facets; ordinary congruence or
equality of touch counts does not by itself provide such congruences.
\end{definition}

\begin{conditionaltheorem}[$n=5$ under BRH]\label{thm:n5}
No partition of a regular tetrahedron into five full-dimensional congruent
convex polytopes satisfies BRH.  In particular, five-piece impossibility is
proved conditionally on BRH, but no unconditional $n=5$ impossibility is
claimed for congruent copies with different boundary roles.
\end{conditionaltheorem}

\begin{proof}
Let $A$ be the area of one face of $T$, and let $t=|B|$.  The boundary
footprints cover the four faces of $T$ up to sets of area zero.  Indeed, two
full-dimensional convex pieces contained in $T$ cannot share a relatively
open subset of an ambient face without overlapping interiors, so their
boundary footprints can overlap only in area-zero sets.  Since BRH identifies
the exposed facets in every copy, their total area is the same in each piece;
hence the common exposed area is
\begin{equation}\label{eq:brh-exposed-area}
 E=\frac{4A}{5}.
\end{equation}
Every exposed facet has area at most $E$, because it is a positive summand of
that total.  Thus every face of $T$ needs at least two boundary footprints.
Convexity gives at most one exposed facet on each ambient face, so $t\le4$.

If $t=0$, there is no boundary coverage.  If $t=1$, there are only five
footprints in total, whereas four faces require at least eight.  Suppose
$t=2$, and write the two facet areas as $a\le b$.  Normalize $A=1$, so
$a+b=4/5$.  There are ten footprints, so at least one face has exactly two
footprints and another has at least three.  A two-footprint face has area
$2a$, $a+b$, or $2b$.  The first two are less than $1$, so $2b=1$ and
$(a,b)=(3/10,1/2)$.  Three footprints then have one of the four sums
$9/10,11/10,13/10,15/10$, while four or more have area at least
$4a=6/5$.  No face with at least three footprints can therefore have area
$1$, a contradiction.

It remains to consider $t=3$ or $t=4$.  Center $T$ at the origin, let
$\nu_0,\ldots,\nu_3$ be its outward unit face normals, and let $h>0$ be
the common support number, so that
\[
 T=\{x\in\mathbb R^3:\nu_j\cdot x\le h\text{ for }j=0,\ldots,3\}.
\]
They satisfy
$\nu_0+\nu_1+\nu_2+\nu_3=0$, and every three are linearly independent.
Replace the reference tile by $P_1$ and set
$\phi_i=\gamma_i\gamma_1^{-1}$ for $i=1,\ldots,5$.  Write
$\phi_i(x)=Q_i x+u_i$.  The map $\phi_i$ sends three exposed facets of
$P_1$ to exposed facets of $P_i$.  At an exposed facet the piece's outward
normal is the \emph{positive} ambient outward normal $+\nu_j$ (not $-\nu_j$),
and isometries transport outward normals.  Thus $Q_i$ sends three distinct tetrahedral face normals to three
distinct tetrahedral face normals.  The zero-sum identity
forces the fourth normal to be sent to the fourth one as well.  Hence $Q_i$
permutes the four normals.

For a matched facet on $\nu_j\cdot x=h$, the image facet lies on the
corresponding ambient support plane.  For $x$ on that facet,
\[
 h=(Q_i\nu_j)\cdot(Q_i x+u_i)=h+(Q_i\nu_j)\cdot u_i.
\]
This makes $u_i$ perpendicular to three independent normals, so $u_i=0$.
Every $\phi_i$ is consequently an element of the full tetrahedral symmetry
group $H$.  Its order is $24$: each permutation of the four vertices of a
regular tetrahedron extends uniquely to an orthogonal symmetry.

For completeness, the required orbit count does not assume that $H$ permutes
the partition.  Choose a point $x$ in the interior of $T$ outside the finite
union of all fixed-point sets of nonidentity elements of $H$ and all inverse
images under $H$ of the piece boundaries.  Such a point exists because this
finite union has three-dimensional measure zero.  Then $O=Hx$ has $24$ points and
meets no piece boundary.  Since $P_i=\phi_i(P_1)$ and $\phi_i\in H$,
$|O\cap P_i|=|O\cap P_1|$ for every $i$.  If this common number is $m$,
partitioning $O$ gives $24=5m$, impossible.  This proves the conditional
theorem.
\end{proof}

The proof does not require the partition to be face-to-face.  Its restriction
is BRH: it does not establish a common boundary-role set for arbitrary
congruent copies.  Establishing that property for a specified class, or
excluding mixed boundary roles directly, remains a separate problem.

\begin{observation}[Exact bookkeeping for a common exposed-area list]
\label{obs:n5-area}
Assume only that each piece has the same list of positive exposed-facet areas
$m=(m_1,\ldots,m_t)$, without requiring the facet identifications to be
induced by congruences.  Boundary coverage again gives $\sum m_r=4A/5$.
Normalize $A=60$, hence $E=48$; this does not require individual areas to be
integers.  Label equal-area occurrences arbitrarily.  Let $N_{jr}$ count
occurrences of label $r$ on ambient face $j$.  Necessary equations are
\[
 N m=60(1,1,1,1)^\mathsf{T},\qquad m_r>0,
\]
with row sums $k_j$ and column sums $5$.  Each $k_j$ lies between $2$ and
$5$, and $\sum k_j=5t$.  For $t=2,3$ the table lists every possible sorted
$k$ and every positive area solution of these count equations.  The script
\path{scripts/n5_real_area_certificate.py} enumerates all $3461$ matrices
with these margins and solves the rational linear systems exactly.  No
consistent rank-deficient system occurs, so no real solution family is lost
by listing isolated rational solutions.
\begin{center}
\begin{tabular}{c c l}
\toprule
$t$ & $k$ & positive sorted area multisets\\
\midrule
2 & $(4,2,2,2)$ & none\\
2 & $(3,3,2,2)$ & none\\
3 & $(5,5,3,2)$ & none\\
3 & $(5,4,4,2)$ & $(6,12,30)$\\
3 & $(5,4,3,3)$ & $(6,15,27)$, $(6,18,24)$, $(48/5,84/5,108/5)$, $(12,12,24)$\\
3 & $(4,4,4,3)$ & $(6,18,24)$, $(12,12,24)$, $(12,16,20)$\\
\bottomrule
\end{tabular}
\end{center}
There are six distinct positive multisets and eight residual occurrences.  The
rational multiset $(48/5,84/5,108/5)$ is genuine at the incidence level; one
aligned witness is
\[
 N=\begin{pmatrix}4&0&1\\1&3&0\\0&1&2\\0&1&2\end{pmatrix},
 \qquad
 m=\left(\frac{48}{5},\frac{84}{5},\frac{108}{5}\right)^\mathsf{T}.
\]
Its row sums are $(5,4,3,3)$ and every weighted row sum is $60$.  The
calculation concerns areas and incidence only; no matrix here is a geometric
dissection certificate.
\end{observation}

\begin{observation}[A local algebraic area model]\label{obs:n5-nn}
The following identity retains a local calculation from the NN first-factor
model, with all algebraic assumptions explicit.  For real $x,y,z,L$ with
$L\ne0$, define
\begin{align*}
 Q_p&=x^2-xy+y^2, & Q_r&=z^2-z+1,\\
 B&=x^2+xy+y^2-(z^2+z+1), & D&=xy-z,\\
 E_\lambda&=Q_r-12L^2Q_p, &
 E_{\rm axis}&=Q_p\bigl(4-(6L-1)^2\bigr)-12B,\\
 A_{\rm floor}&=D+2LQ_p+\frac{Q_r}{6L}, &
 A_{\rm total}&=A_{\rm floor}+B+D.
\end{align*}
Then the identity
\[
 24L\left(A_{\rm floor}-\frac34A_{\rm total}\right)
 =2L E_{\rm axis}+(1-6L)E_\lambda
\]
holds identically; substituting $B=Q_p-Q_r+2D$ verifies it directly.
If $E_\lambda=E_{\rm axis}=0$ and $A_{\rm total}>0$, it follows that
$A_{\rm floor}/A_{\rm total}=3/4$.  The identity is checked symbolically by
\path{scripts/n5_nn_identity_certificate.py}.
If these quantities describe one exposed facet and the total exposed area
of a piece in the $t=3$ common-area model of \cref{obs:n5-area}, they are
incompatible with its maximum facet-area ratio $30/48=5/8$.
This last conclusion requires that geometric identification.  No assertion
that arbitrary pieces, all NN branches, or mixed boundary roles enter this
algebraic model is made.
\end{observation}

\begin{remark}[Withdrawn metric shortcuts]
The earlier $f=5$ triangular-prism exclusions are not used here.  A triangular-
prism face lattice does not imply parallelism of the two triangular faces.
In affine coordinates with $T=\{x_i\ge0,\ \sum_i x_i\le1\}$, take
$0<a<\min_i b_i$ and $\max_i b_i\le1$.  The polytope
\[
 x_i\ge0,\qquad x_1+x_2+x_3\ge a,\qquad
 \sum_i x_i/b_i\le1
\]
has vertices $ae_i,b_ie_i$ and the triangular-prism face lattice; its two
triangular facets are nonparallel when the $b_i$ are not all equal.
Likewise, a setwise facet correspondence does
not by itself imply pointwise fixing or an orientation-reversing reflection.
The old $H$-cocycle closure is therefore not imported into the present theorem.
\end{remark}

\begin{remark}[Current $n=5$ boundary]
The active result is exactly \Cref{thm:n5}, the conditional BRH obstruction,
together with the finite area and local-chart observations above.  The
unrestricted convex face-to-face five-piece question remains open.  The
historical phase-26 scripts and result files are retained for auditability but
are no longer an active proof chain.  The reproducer runs the BRH arithmetic
checks, the area enumeration, the local algebraic identity, and the auxiliary
atlas checks described in the repository README.  The BRH proof is a written
geometric argument, not an inference from passing test counts.
\end{remark}

\subsection{Summary of the programme}\label{ssec:prog-summary}

\cref{conj:transitivity} is equivalent to \cref{conj:divisor} combined with
\Cref{cor:forward}.  Its current evidence level is \textsf{L1+}, supported
by
\begin{itemize}
  \item one \textsf{L2} theorem (\cref{lem:stab-transport}, Stabiliser Transport);
  \item one \textsf{L1+} proposition (\cref{prop:csl}, Conditional Separation,
        conditional on Hypothesis H);
  \item one conjecture at \textsf{L1} (\cref{conj:bn-orbits}, Boundary-Normal
        Orbits);
  \item a Spherical Link Program constraining pieces-per-vertex to
        $k \in \{1, 3, 6, \ldots\}$ within a uniform convex model under
        Hypothesis H;
  \item the dihedral rigidity programme (\Cref{ssec:rigidity}) reducing
        \cref{conj:transitivity} to GAPs A, B;
  \item the $n = 16$ all-one-face conditional reduction
        (\cref{thm:n16-aof-coxeter,thm:n16-aof-conditional}), \textsf{L3} in the
        Coxeter-pure case and \textsf{L2} (conditional on
        $(\Hroles)+(\Hslot)$) in the general convex case;
  \item and the conditional $n = 5$ BRH obstruction
        (\cref{thm:n5}).  The unrestricted case with different boundary roles
        remains open; exact area enumeration and the local NN identity provide
        diagnostics but do not close that gap.
\end{itemize}

%%%%%%%%%%%%%%%%%%%%%%%%%%%%%%%%%%%%%%%%%%%%%%%%%%%%%%%%%%%%%%%%%%%%%%%%
\section{Further questions}\label{sec:further}

\begin{question}[Hinged dissections]\label{q:hinged}
Can any of the eight known dissections be realised as a physical hinged
mechanism?  Edge-hinges ($1$-DOF) are natural candidates for $S_{2}$ and
$S_{4}$, vertex-hinges ($3$-DOF) for the dense configurations $S_{24}$.
No results are known in this direction.  Demaine's universality theorems
on scissors-congruent polyhedra~\cite{DemaineDemaine2007} apply to pairs
of polyhedra, not to partitions of a single polyhedron; their adaptation
to the present setting is open.
\end{question}

\begin{question}[Higher-dimensional extension]\label{q:higher-dim}
For the regular $d$-simplex with symmetry group $S_{d+1}$ of order
$(d + 1)!$, does $n$ admit a congruent dissection if and only if
$n \mid (d + 1)!$?  The answer is \emph{no} for $d = 2$: the equilateral
triangle admits congruent dissection into any $n = k^{2}$ pieces via
self-similar subdivision.  The tetrahedron ($d = 3$) may be the critical
dimension at which the Dehn invariant first constrains the admissible
$n$; this should be a theorem for $d = 3$, with modifications for
$d \geq 4$.
\end{question}

\begin{question}[Burnside ring extensions]\label{q:burnside-ext}
The Burnside ring framework of \Cref{sec:burnside} opens questions in
\emph{equivariant $K$-theory} of simplicial complexes and connections to
the representation ring $R(\Td)$.  The relationship between the mark
table and the character table of $S_{4}$ may encode deeper
number-theoretic structure.
\end{question}

%%%%%%%%%%%%%%%%%%%%%%%%%%%%%%%%%%%%%%%%%%%%%%%%%%%%%%%%%%%%%%%%%%%%%%%%
\appendix

\section{Computational reproduction sketch}\label{app:computational}

The released metric, Dehn, stabiliser, and SAT-style certificates use exact
rational arithmetic, including Python's built-in \texttt{Fractions} module.
We include below a minimal self-contained reference implementation for the
Dehn-invariant verification at $n = 24$ (orthoscheme).  The public replay suite
for \cref{thm:n5}, together with the auxiliary atlas checks, is available in
\path{scripts/} and \path{results/} of the companion repository identified in
\Cref{sec:intro}.  The active $n=5$ script certifies only the finite arithmetic
checks used by the BRH proof; historical phase-26 reduction scripts are not
part of that active chain.  As stated there, the two general-convex hypotheses of
\cref{thm:n16-aof-conditional} are not computationally discharged by this
public package.

\Needspace*{22\baselineskip}
\begin{lstlisting}[language=Python,caption={Minimal SymPy verification of
the Dehn scaling $D(S_{24}) = D(T)/24$ for the orthoscheme piece.},
label={lst:dehn-24}]
import sympy as sp
from sympy import Rational, Point3D, sqrt
# Regular tetrahedron T with edge length 1
A = Point3D(0, 0, 0)
B = Point3D(1, 0, 0)
C = Point3D(Rational(1, 2), sqrt(3)/2, 0)
D = Point3D(Rational(1, 2), sqrt(3)/6, sqrt(6)/3)
# Canonical orthoscheme piece S_24
M_AB   = (A + B) / 2
G_ABC  = (A + B + C) / 3
G_T    = (A + B + C + D) / 4
S24    = [A, M_AB, G_ABC, G_T]
# Edge enumeration, dihedral-angle computation via dot products of
# outward face normals, and Dehn tensor product in R (x)_Z (R/pi Q)
# follow directly; total runtime < 1 s.
\end{lstlisting}

\section{Evidence-level glossary}\label{app:evidence}\label{sec:evidence-glossary}

\begin{description}
\item[\textsf{L0}]
  \emph{Heuristic.}  A claim supported by plausibility or partial checks
  but without a formal proof attempt.
\item[\textsf{L1}]
  \emph{Partial evidence.}  A claim supported by substantial evidence
  (e.g.\ computational enumeration up to a fixed bound) but lacking a
  closed proof.
\item[\textsf{L1+}]
  \emph{Strong partial evidence.}  A claim supported by a partial proof
  together with the explicit identification of the remaining residual
  gap(s); typically a reduction to a named sub-lemma.
\item[\textsf{L2}]
  \emph{Conditional theorem.}  A rigorously proved implication whose
  conclusion depends on an explicit, named auxiliary hypothesis (e.g.\
  convexity, face-to-face, or a named hypothesis H).
\item[\textsf{L2+}]
  \emph{Finite-check conditional theorem.}  Like \textsf{L2}, but each
  standing auxiliary hypothesis is a finite, explicitly verifiable
  condition on a candidate configuration (e.g.\ a cocycle identity,
  an orbit table, or a single-piece combinatorial type exclusion).
\item[\textsf{L3}]
  \emph{Unconditional theorem.}  A fully proved statement with no
  auxiliary hypothesis beyond those of the ambient setting (convex
  face-to-face congruent dissection).
\end{description}

%%%%%%%%%%%%%%%%%%%%%%%%%%%%%%%%%%%%%%%%%%%%%%%%%%%%%%%%%%%%%%%%%%%%%%%%
\section{The \texorpdfstring{$24$}{24}-cell and quaternionic action}\label{app:24cell}

This appendix provides a \emph{secondary} algebraic viewpoint.  The core
results of the paper (\Cref{sec:atlas}--\Cref{sec:burnside},
\Cref{sec:ifs}--\Cref{sec:nontrans}) do not depend on it; it serves as a
partial independent verification of the dissection inventory through the
quaternionic / $24$-cell correspondence, and may be skipped on first
reading.

\subsection{The binary tetrahedral group \texorpdfstring{$2T$}{2T}}

The binary tetrahedral group $2T \subset S^{3} \subset \HH$ consists of $24$
unit quaternions: the $8$ \emph{type-$1$} units $\{\pm 1, \pm i, \pm j, \pm k\}$
and the $16$ \emph{type-$2$} units $\tfrac{1}{2}(\pm 1 \pm i \pm j \pm k)$.
These are precisely the $24$ \emph{Hurwitz units}, the vertices of the
regular $24$-cell inscribed in $S^{3}$.  The group has seven conjugacy
classes of sizes $(1,6,1,4,4,4,4)$.  Its centre is $\{\pm 1\}$, and
\begin{equation}\label{eq:2T-ses}
  1 \longrightarrow \{\pm 1\} \longrightarrow 2T \xrightarrow{\;\pi\;}
  A_{4} \longrightarrow 1
\end{equation}
is exact.

\subsection{Triple bijection}

The quaternionic conjugation action $v \mapsto q v \overline{q}$ on pure
quaternions $v$ yields a canonical triple bijection
\begin{equation}\label{eq:triple-bij}
  \underbrace{24 \text{ Hurwitz units}}_{2T}
  \;\longleftrightarrow\;
  \underbrace{24 \text{ permutations}}_{\Td \cong S_{4}}
  \;\longleftrightarrow\;
  \underbrace{24 \text{ orthoschemes}}_{\text{flags}(T)}.
\end{equation}
Proper rotations give $24 \twoheadrightarrow 12$ even permutations ($A_{4}$)
with kernel $\{\pm 1\}$; composition with a fixed involution $\sigma$
recovers the $12$ odd permutations ($S_{4} \setminus A_{4}$); and each of the
$24$ orthoschemes in the barycentric subdivision of $T$ is hit exactly once.

\subsection{Subgroup lattice and the \texorpdfstring{$n = 2$}{n = 2} gap}

The subgroup lattice of $2T$ maps under $\pi$ to the subgroup lattice of
$A_{4} \leq S_{4}$.  Seven of the eight dissection numbers are recovered
directly from coset spaces $2T/H$: $\{1,3,4,6,8,12,24\}$
(\cref{tab:2T-lattice}).  The missing case $n = 2$ arises because $2T$ has
no subgroup of order $12$; indeed $\pi$ has image $A_{4} \leq S_{4}$ of
index $2$, and all $|\pi(H)|$ divide $12$.  The $n = 2$ dissection is
captured by the $S_{4}$-based Burnside analysis (\Cref{sec:burnside}) via
stabiliser $A_{4}$; see \cref{tab:2T-lattice}.

\begin{table}[!t]
\centering\small
\caption{Subgroups of $2T$ and dissection matches.}
\label{tab:2T-lattice}
\begin{tabular}{@{} l c c l c c @{}}
\toprule
$2T$-subgroup & $|H|$ & $[2T:H]$ & $\pi(H) \leq A_{4}$ & $|\pi(H)|$ & Dissection \\
\midrule
$\{e\}$             & $1$  & $24$ & $\{e\}$      & $1$ & $S_{24}$ \\
$\Z_{2}$ (centre)    & $2$  & $12$ & $\{e\}$      & $1$ & $S_{12}$ \\
$\Z_{3}$            & $3$  & $8$  & $\Z_{3}$     & $3$ & $S_{8}$  \\
$\Z_{4}$            & $4$  & $6$  & $\Z_{2}$     & $2$ & $S_{6}$  \\
$\Z_{6}$            & $6$  & $4$  & $\Z_{3}$     & $3$ & $S_{4}$  \\
$Q_{8}$             & $8$  & $3$  & $V_{4}$      & $4$ & $S_{3}$  \\
$2T$                & $24$ & $1$  & $A_{4}$      & $12$& $S_{1}$  \\
\bottomrule
\end{tabular}
\end{table}

%%%%%%%%%%%%%%%%%%%%%%%%%%%%%%%%%%%%%%%%%%%%%%%%%%%%%%%%%%%%%%%%%%%%%%%%
\begin{thebibliography}{99}

\bibitem{AW1971}
G.~L. Alexanderson and J.~E. Wetzel,
\emph{Dissections of a tetrahedron},
J.\ Combinatorial Theory Ser.\ A \textbf{11} (1971), 58--66.

\bibitem{Burnside1911}
W.~Burnside,
\emph{Theory of Groups of Finite Order},
second ed., Cambridge University Press, 1911; reprinted Dover, 1955.

\bibitem{ConwayJones1976}
J.~H. Conway and A.~J. Jones,
\emph{Trigonometric diophantine equations (On vanishing sums of roots of unity)},
Acta Arith.\ \textbf{30} (1976), 229--240.

\bibitem{CoxeterPure2026}
O.~Babanskyy,
\emph{Congruent Coxeter-pure partitions of the regular tetrahedron:
catalogue, convexity, and classification},
companion paper, preprint, 2026.

\bibitem{DemaineDemaine2007}
E.~D.\ Demaine and M.~L.\ Demaine,
\emph{Jigsaw puzzles, edge matching, and polyomino packing: connections and
complexity},
Graphs Combin.\ \textbf{23} (2007), 195--208.

\bibitem{Dupont2001}
J.~L.\ Dupont,
\emph{Scissors congruences, group homology and characteristic classes},
Nankai Tracts in Mathematics, vol.~1, World Scientific, 2001.

\bibitem{Hadwiger1951}
H.~Hadwiger,
\emph{Zerlegungsgleichheit und additive Polyederfunktionale},
Arch.\ Math.\ \textbf{2} (1951), 307--310.

\bibitem{Hill1895}
M.~J.~M. Hill,
\emph{Determination of the volumes of certain species of tetrahedra
without employment of the method of limits},
Proc.\ London Math.\ Soc.\ (1) \textbf{27} (1895), 39--53.

\bibitem{Hiller1985}
H.~Hiller,
\emph{The crystallographic restriction in higher dimensions},
Acta Crystallogr.\ Sect.\ A \textbf{41} (1985), 541--544.

\bibitem{Ong1991}
H.~L.\ Ong,
\emph{A bisection algorithm for tetrahedral subdivision},
in: Computing in Euclidean Geometry (D.--Z.\ Du and F.\ Hwang, eds.),
World Scientific, 1992, pp.~327--347.

\bibitem{PPS2023}
M.~A.\ Padr\'on, \'A.\ Plaza and J.~P.\ Su\'arez,
\emph{Congruence classes of tetrahedra generated by longest-edge bisection},
Appl.\ Math.\ Comput.\ \textbf{441} (2023), Paper No.\ 127707, 10 pp.

\bibitem{Sydler1965}
J.-P.\ Sydler,
\emph{Conditions n\'ecessaires et suffisantes pour l'\'equivalence des
poly\`edres de l'espace euclidien \`a trois dimensions},
Comment.\ Math.\ Helv.\ \textbf{40} (1965), 43--80.

\bibitem{tDieck1979}
T.~tom Dieck,
\emph{Transformation groups and representation theory},
Lecture Notes in Mathematics, vol.~766, Springer-Verlag, 1979.

\end{thebibliography}

\end{document}
